\documentclass[11pt,reqno]{amsart}
\usepackage[T1]{fontenc}
\usepackage{lmodern,amsmath,amssymb,amsthm,mathtools,mathrsfs}
\usepackage[a4paper,margin=27mm]{geometry}
\usepackage{booktabs,longtable,array,microtype,needspace,etoolbox}
\usepackage[dvipsnames]{xcolor}
\definecolor{InternalLink}{HTML}{0057B8}
\definecolor{CitationLink}{HTML}{1B7837}
\definecolor{URLLink}{HTML}{A33A00}
\definecolor{ReturnLink}{HTML}{6A3D9A}
\usepackage{hyperref}
\hypersetup{colorlinks=true,linkcolor=InternalLink,citecolor=CitationLink,
urlcolor=URLLink,linktoc=all,pdfborder={0 0 0},
pdftitle={Locally bounded classical plurisubharmonic functions and mixed currents},
pdfauthor={}}
\usepackage[nameinlink,noabbrev]{cleveref}
\newtheorem{theorem}{Theorem}[section]
\newtheorem{proposition}[theorem]{Proposition}
\newtheorem{lemma}[theorem]{Lemma}
\newtheorem{corollary}[theorem]{Corollary}
\theoremstyle{definition}
\newtheorem{definition}[theorem]{Definition}
\newtheorem{convention}[theorem]{Convention}
\theoremstyle{remark}
\newtheorem{remark}[theorem]{Remark}
\numberwithin{equation}{section}
\newcommand{\an}{\mathrm{an}}
\newcommand{\dd}{\mathrm d'\mathrm d''}
\newcommand{\dpr}{\mathrm d'}
\newcommand{\dpp}{\mathrm d''}
\newcommand{\cO}{\mathcal O}
\newcommand{\Cb}{\mathcal C_b}
\newcommand{\Rb}{\mathcal R_b}
\newcommand{\PLsp}{\mathrm{PL}_{\mathrm{sp}}}
\newcommand{\MA}{\operatorname{MA}}
\newcommand{\supp}{\operatorname{supp}}
\newcommand{\returnlink}[2]{\par\smallskip\noindent
\hyperref[#1]{\textcolor{ReturnLink}{Return to #2.}}\par}
\newcommand{\KeepStatement}[1]{\BeforeBeginEnvironment{#1}{\Needspace{5\baselineskip}}}
\forcsvlist{\KeepStatement}{theorem,proposition,lemma,corollary,definition,convention,remark}
\AddToHook{cmd/section/before}{\Needspace{8\baselineskip}}
\AddToHook{cmd/subsection/before}{\Needspace{6\baselineskip}}
\crefname{convention}{convention}{conventions}
\Crefname{convention}{Convention}{Conventions}
\setlength{\emergencystretch}{2em}
\title[Classical psh functions and mixed currents]{Locally bounded classical plurisubharmonic functions\newline and mixed currents}
\author{}
\date{10 November 2026}

\begin{document}
\begin{abstract}
Let $K$ be a finite extension of $\mathbb Q_p$ and let $X$ be a good,
Hausdorff, boundaryless, quasi-smooth $K$-analytic space. We prove that
every locally bounded classically plurisubharmonic function on $X$
is locally a decreasing limit of semipositive rational piecewise-linear
functions defined over $K$ on one fixed neighborhood. We use this
regularization to construct mixed currents in the original
Chambert--Loir--Ducros calculus. The unweighted currents are symmetric,
closed and strongly positive. Their coefficient measures are Radon
measures on the analytic space, so multiplication by a locally bounded
classical weight uses its given Borel representative. These weighted
products are jointly continuous along decreasing sequences with locally
bounded limits, and the construction is compatible with \'etale
pullback of the potentials. The regularization proof uses curve slicing,
projective contact rigidity, and continuity of ample envelopes on
normal models. A fixed vertical correction for mixed-characteristic
test ideals provides the envelope estimate.
\end{abstract}
\maketitle

\section{Introduction}\label{sec:introduction}
Let $K$ be a finite extension of $\mathbb Q_p$. A classically
plurisubharmonic function on a Berkovich $K$-analytic space is an
upper-semicontinuous function whose pullback to every analytic curve,
after any complete extension of $K$, is subharmonic \cite{GR}.
This definition tests positivity on curves and does not require the
function to have a smooth local expression. The question addressed here
is whether locally bounded functions in this class admit decreasing
piecewise-linear approximation, and whether their curvatures have mixed
products in the Chambert--Loir--Ducros calculus.

For continuous semipositive potentials, uniform approximation provides
the mixed products. A locally bounded classical potential may be
discontinuous. The recursive construction then requires multiplication
of a current by a Borel function, and its continuity requires a limit
theorem in which both that function and the current vary. We prove
these assertions for the original CLD test forms. The approximation
takes place over $K$ on a fixed neighborhood of each point; all weighted
products use the given pointwise representatives.

\subsection{Definitions and main results}\label{sec:main-results}

\begin{convention}\label[convention]{conv:spaces}
Fix a finite extension $K/\mathbb Q_p$ with its prescribed absolute
value. Choose a uniformizer $\pi$ and put $\lambda=-\log|\pi|$.
Unless another field is specified, $X$ is a good, Hausdorff,
boundaryless, quasi-smooth $K$-analytic space of pure dimension
$n\geq1$. All opens are ordinary topological opens. A real function
is \emph{locally bounded} if it is bounded above and below on a
neighborhood of every point. No replacement on a null or polar set
is made.
\end{convention}

\begin{definition}\label[definition]{def:classes}
For an ordinary open $U\subset X$, let $\Cb(U)$ be the real-valued, locally bounded,
upper-semicontinuous functions $u$ such that, for every complete
isometric extension $F/K$ and every morphism $C\to U_F$ from an
$F$-analytic curve, the pullback of $u$ is subharmonic in the sense
of \cite[Definition~4.1 and \S3]{GR}.

Let $\PLsp(U)$ denote the continuous real piecewise-linear functions
semipositive in the local model, or equivalently reduction, sense of
\cite[Theorem~4.3]{GR}. Define $\Rb(U)$ to be the locally
bounded functions $u$ such that every point has one ordinary
neighborhood $V\subset U$ and a sequence
\[
 p_m\in\PLsp(V),\qquad p_m\downarrow u|_V.
\]
The neighborhood is independent of $m$. On an algebraic chart,
\emph{rational} PL means rational slopes and constants in
$\lambda\mathbb Q$.
\end{definition}

\begin{convention}\label[convention]{conv:forms}
We use the original CLD forms $A^{p,q}$ and their compactly
supported test spaces $A_c^{p,q}$, with the canonical calibration.
Put $D=\dd$. A current of bidegree $(k,k)$ acts on
$A_c^{n-k,n-k}$. A symmetric current is called \emph{strongly
positive} if it is nonnegative on weakly positive complementary
original forms. The weakly positive cone is tested by decomposable
vectors; the strongly positive cone is generated by decomposable
squares. These conventions concern the original test sheaf, not the
larger weakly smooth sheaf of \cite{GJR}.

A metric is written additively in weights $-\log\|e\|$, so its
curvature is $D[-\log\|e\|]$. On a projective variety a continuous
\emph{model-semipositive} metric means a uniform limit of nef
rational model metrics on the same identified line bundle.
A \emph{bounded classical metric} is a metric whose local weights
are classical and whose difference from a continuous reference
metric is bounded. A rational \emph{pure Fubini--Study metric},
abbreviated pure-FS, is given by a finite generating tuple
$s_0,\ldots,s_N$ of a tensor power $mL$, with local weight
\[
 \frac1m\max_i\bigl(\log|s_i/e^{\otimes m}|+c_i\bigr),
 \qquad c_i\in\lambda\mathbb Q .
\]
After replacing $m$ by a multiple, the constants can be absorbed
by multiplying the sections by powers of $\pi$.
\end{convention}

The recursive formula for mixed currents will involve multiplying a
current by a locally bounded, possibly discontinuous function. The
following condition specifies the measures used for that multiplication.

\begin{definition}[Coefficient Radon currents]\label[definition]{def:coefficient-radon}
A current $T$ is of \emph{coefficient Radon type} if each finite
original moment presentation represents its coefficient functionals
by signed Radon measures on the analytic domain, compatibly
with all original smooth scalar functions, relations between forms,
and common refinements. In particular, for every compact original
test $\eta$, there is a signed Radon measure $\mu_{T,\eta}$ with
\[
 \int a\,d\mu_{T,\eta}=\langle T,a\eta\rangle
 \quad\text{for every original smooth scalar }a.
\]
For a locally bounded Borel function $h$, multiplication is
defined by
\begin{equation}\label{eq:borel-multiplication}
 \langle hT,\eta\rangle=\int h\,d\mu_{T,\eta}.
\end{equation}
The compatibility requirements make this definition independent of
the chosen presentation. A measure only on one tropical image does
not suffice for this definition.
\end{definition}

\begin{theorem}[Local decreasing regularization]
\label[theorem]{thm:local-regularization}
For every ordinary open $U\subset X$,
\[
 \Cb(U)=\Rb(U).
\]
For $u\in\Cb(U)$, the local approximants may be chosen rational PL,
defined over $K$ on one fixed neighborhood. Both classes are stable
under decreasing sequences whose pointwise limit is locally
bounded.
\end{theorem}
\noindent\hyperref[proof:local-regularization]{Proof on
p.~\pageref{proof:local-regularization}.}

\Needspace{15\baselineskip}
\begin{theorem}[Mixed currents]\label[theorem]{thm:mixed-currents}
For $u_1,\ldots,u_k\in\Cb(U)$, $0\leq k\leq n$, there are
coefficient Radon currents
\begin{equation}\label{eq:recursive-current}
 T_0=[1],\qquad
 T_k(u_1,\ldots,u_k)=D\bigl(u_1T_{k-1}(u_2,\ldots,u_k)\bigr).
\end{equation}
They are symmetric, $\dpr$- and $\dpp$-closed, strongly positive,
and multilinear on the cone in the differentiated arguments. They agree
with the continuous CLD products whenever the potentials are
continuous and semipositive. They are independent of decreasing PL
regularization and commute with ordinary restriction and gluing.
For $u_0\in\Cb(U)$, the weighted current
\begin{equation}\label{eq:source-current}
 S_k(u_0;u_1,\ldots,u_k)=u_0T_k(u_1,\ldots,u_k)
\end{equation}
uses the Borel multiplication of
\Cref{def:coefficient-radon}.
\end{theorem}
\noindent\hyperref[proof:main-currents]{Proof on
p.~\pageref{proof:main-currents}.}

\begin{theorem}[Joint decreasing continuity]
\label[theorem]{thm:joint-monotone}
Suppose that, for $0\leq i\leq k$, the functions
$u_{i,j}\in\Cb(U)$ decrease pointwise to a locally
bounded function $u_i$. Then
\[
 S_k(u_{0,j};u_{1,j},\ldots,u_{k,j})
 \longrightarrow S_k(u_0;u_1,\ldots,u_k)
\]
on every form in $A_c^{n-k,n-k}(U)$. The corresponding unweighted
currents also converge. In top degree these are locally weak
convergences of Radon measures, against every continuous
compactly supported scalar function.
\end{theorem}
\noindent\hyperref[proof:main-currents]{Proof on
p.~\pageref{proof:main-currents}.}

\begin{corollary}\label[corollary]{cor:quasiprojective}
The preceding results apply on every ordinary open in the
analytification of a smooth quasi-projective $K$-variety. Local
boundedness there imposes no boundedness condition at a deleted
algebraic boundary.
\end{corollary}

\begin{remark}[The undifferentiated argument]\label[remark]{rem:source-slot}
A constant in a differentiated argument has zero curvature. A change
$u_0\mapsto u_0+c$ in the undifferentiated argument changes $S_k$ by $cT_k$.
Thus $S_k$ is symmetric in $u_1,\ldots,u_k$, but not in all its
arguments. The weighted current may be signed and need not be closed.
The hypothesis on the limit in \Cref{thm:joint-monotone} is
essential to the stated class; for example, the decreasing constants
$-j$ have no locally bounded limit.
\end{remark}

\subsection{Proof strategy and relation to earlier work}
The proof has two parts. We first pass from the curve condition to
local decreasing regularization. We then use this regularization to
construct the mixed currents and prove their continuity. The first
part uses only first-order currents and products of continuous metrics,
so it does not presuppose the products in \Cref{thm:mixed-currents}.

The continuous products come from uniform approximation by semipositive
model metrics. Approximation by Fubini--Study metrics and by smooth
semipositive metrics is supplied by
\cite[Corollary~7.6 and Theorems~7.8, 7.14]{BE}, in the current calculus
of \cite{CLD12}. To treat a discontinuous classical potential $v$, we
slice a test presentation by algebraic curves. The curve condition
applies over the completed residue fields of the slicing points, even
when those fields are nondiscretely valued. Thuillier's lisse
approximation and Dirichlet energy \cite[\S\S3.4--3.5]{Thuillier}
give a positive Laplacian on each curve. Integrating these Laplacians
gives the first-order current $D[v]$, its coefficient measures, and the
Green formulas needed below. The lisse functions are auxiliary
piecewise-affine potentials on skeleta; the tests remain original CLD
forms. We use the skeleton integration formulas of \cite{CLD25} on
their original unit-moment presentations.

The passage to regularization is made on a smooth projective
geometrically integral variety over a field in the finite unit-Gauss
tower class $\mathcal G_K$ of \Cref{def:gauss-fields}, with an ample
line bundle $L$. The contact theorem states that a bounded
classical metric $r$ and a continuous model-semipositive metric $p$
satisfy
\[
 r\geq p,\qquad \int (r-p)c_1(L,p)^n=0
 \quad\Longrightarrow\quad r=p
\]
pointwise. Its proof is an induction on dimension. The first-order
Green estimate controls the gap $r-p$ against smooth backgrounds.
A generic hyperplane then identifies its limit against the continuous
background and reduces the vanishing assertion to dimension $n-1$.
The hyperplane is defined over a unit-Gauss completion, which is why
the induction treats arbitrary data over all finite unit-Gauss towers
at once.

To obtain the approximants, choose continuous metrics
$\rho_\ell\downarrow r$ and take their semipositive envelopes
$s_\ell=P_L\rho_\ell$. Continuity and orthogonality of these envelopes
give
\[
 \int (\rho_\ell-s_\ell)c_1(L,s_\ell)^n=0.
\]
Apply contact to $\max(r,s_\ell)$ and $s_\ell$. Their nonnegative gap
is bounded by $\rho_\ell-s_\ell$, so contact gives
$r\leq s_\ell\leq\rho_\ell$ and hence $s_\ell\downarrow r$.
The envelope theorem is proved in \Cref{sec:normal-envelopes} on a
normal model. Its algebraic input is the corrected subadditivity
inequality of \Cref{thm:fixed-correction}; the vertical correction is
independent of the ideals and their exponents and therefore contributes
an error tending to zero after normalization by degree. To return to
a local potential, \Cref{lem:common-collar} realizes it, up to a
constant, as a metric weight on a smaller fixed neighborhood.

For the mixed-current construction, the induction is on exterior
degree. Strong positivity supplies coefficient Radon measures on the
analytic space, making $u_1T_{k-1}$ meaningful before we differentiate
it. To prove weighted convergence, we use Wu's top-degree theorem for
decreasing piecewise-linear potentials
\cite[Theorem~4.2(4)]{Wu}. A complementary test is reduced to products
of curvatures of auxiliary quadratic potentials. In middle degrees,
this reduction is only modulo the annihilator of decomposable squares.
Strong positivity kills that term in the current pairing, and the
corresponding coefficient measure remains zero after Borel
multiplication. Thus the reduction applies to every original
complementary test.

It remains to replace decreasing PL sequences by arbitrary decreasing
classical potentials. We first place the whole family on one projective
variety, using the common bounds supplied by its first terms and its
limits. A compact diagonal argument then chooses decreasing PL
approximants whose current pairings approach the required ones. Choosing
the common local domain before the approximants resolves the domain
dependence noted in \cite[Remark~4.6]{Wu} for this class of functions.

\Cref{sec:first-order,sec:contact} establish first-order positivity and
regularization; \Cref{sec:construction,sec:arbitrary-columns} construct
the products and prove decreasing continuity. \Cref{sec:localization}
transports these results to ordinary analytic charts and glues the
currents. \Cref{sec:etale} treats specified potential tuples under
\'etale pullback, including the finite-degree formula. The curve
arguments, normal-model envelope construction, and common local domain
are proved in the appendices. The
\hyperref[sec:notation]{index of notation} records the definitions and
normalizations used throughout.

\enlargethispage{\baselineskip}
\begingroup
\footnotesize
\noindent\begin{minipage}{\textwidth}
\tableofcontents
\end{minipage}
\endgroup

\section{First-order currents and energy}\label{sec:first-order}

We construct $D[v]$ by integrating curve Laplacians in the directions
of complementary test forms. The resulting coefficient measures allow
us to form the quadratic energy expressions for bounded potentials.
These first-order expressions supply the contact estimate in
\Cref{thm:projective-regularization}; no higher product of discontinuous
potentials is used at this stage.

\begin{definition}\label[definition]{def:gauss-fields}
Let $\mathcal G_K$ be the finite unit-Gauss tower class of
\Cref{lem:gauss-countable}. The induction below is carried out
simultaneously over every $F\in\mathcal G_K$, for arbitrary data
defined over $F$. The residue field of such an $F$ need not be finite
or perfect. Its valuation ring is nevertheless a complete excellent
mixed-characteristic DVR, and its compact analytic domains used here
are metrizable.
\end{definition}

For a current of bidegree $(1,1)$, the complementary degree is $n-1$.
Every $(n-1)$-vector on a face of dimension at most $n$
is decomposable. This is the linear-algebra fact that permits the
following reduction to curve-slicing directions.

\begin{lemma}[One-curvature complementary spanning]
\label[lemma]{lem:first-order-span}
Let $r=n-1$. On a finite original moment presentation, every
symmetric compact $(r,r)$-test is, after finite localization, a
signed sum
\begin{equation}\label{eq:first-order-span}
 \eta=\sum_\ell a_\ell\Omega_\ell,\qquad
 \Omega_\ell=\bigwedge_{b=1}^{r}\dpr\ell_{\ell b}\wedge\dpp\ell_{\ell b},
\end{equation}
where the $a_\ell$ are smooth compact scalars and each
$\ell_{\ell b}$ is a rational linear combination of logarithmic
norms of units.
\end{lemma}
\begin{proof}
Let $E$ be the ambient moment tangent space. The orthogonal
complement of the span of decomposable squares in
$\operatorname{Sym}^2(\bigwedge^r E^*)$ consists of quadratic forms
vanishing on decomposable $r$-vectors. Every $r$-vector in a face
tangent space of dimension at most $r+1=n$ is decomposable.
Polarization therefore makes this complement zero on every
face. Rational decomposable squares span the remaining finite
dimensional space. Choose a finite basis from them and apply a
fixed linear splitting to the smooth coefficient matrix of $\eta$.
The scalar coefficients remain smooth. Clear rational denominators
and use a finite original smooth partition of unity.
For $n=1$ the assertion is the scalar statement with empty wedge.
\end{proof}

\begin{proposition}[First-order measures and positivity]
\label[proposition]{prop:first-order}
Let $F\in\mathcal G_K$, let $X/F$ be smooth algebraic of pure
dimension $n$, and let $v$ be locally bounded classical psh on an
ordinary open $U\subset X^{\an}$. Then $D[v]$ is a closed positive
original current of coefficient Radon type. In every algebraic-unit
curve-slicing presentation it satisfies the Borel kernel
identity of \Cref{thm:actual-kernel}.
\end{proposition}
\begin{proof}
By \Cref{lem:exact-units}, the finitely many analytic units in a
localized test presentation may be replaced, with exactly the same
norms, by algebraic units on a fixed smaller norm-box interior.
After clearing denominators, each direction in
\eqref{eq:first-order-span} is given by $n-1$ algebraic units.
These units define a map $p:W\to\mathbb G_m^{n-1}$. When this map
is dominant, its fibers over the base skeleton are the curves on
which we use the subharmonicity of $v$. Choose a nonnegative coordinate top
form on the base equal to its constant volume form near the
projected test support. If $p$ is dominant,
\Cref{thm:actual-kernel} shows that
$a\mapsto\langle D[v],a\Omega_\ell\rangle$ is a positive Radon
measure on $W$. If $p$ is not dominant, the pairing
vanishes by tropical dimension. The signed spanning lemma gives
order-zero bounds in every coefficient direction. Smooth scalar
relations and common refinements follow from Radon uniqueness.

The slicing argument gives positivity in constant decomposable
directions. To obtain positivity on every complementary test, we must
also show that coefficient mass on a face has tangent indices on that
face. Closedness supplies this tangency. Multiply the coefficient measures
by a nonnegative compact smooth cutoff and push them to moment
space. Write their symmetric matrix as $M_{IJ}$, and view each
column as an $r$-vector measure $F_J$. Closedness of $D[v]$ implies
that its distributional boundary is a measure; localization adds
only the derivative of the cutoff times the order-zero
current. For an affine linear equation $\ell$, choose
$\ell_\epsilon=\ell\theta(\ell/\epsilon)$, where $\theta=1$ near
zero and has compact support. This cutoff tends uniformly to zero,
while its differential detects precisely the mass on $\{\ell=0\}$. Insert $\ell_\epsilon$ times a compact
scalar and a constant $(r-1)$-covector into the boundary equation.
Terms containing $\ell_\epsilon$ tend to zero. Its differential is
bounded and converges pointwise to
$\mathbf1_{\{\ell=0\}}\,d\ell$. Dominated convergence for the
finite coefficient measures gives
\[
 (\iota_{d\ell}F_J)|_{\{\ell=0\}}=0.
\]
The equations of the affine span of each face therefore make both
indices tangent to that face. On its tangent space every $r$-vector
is decomposable. Positivity on rational decomposable directions,
followed by density, makes the coefficient matrix positive
semidefinite there. It pairs nonnegatively with every positive
original test. For $n=1$ positivity is the curve statement.

The coefficient identities are identities of Radon measures and
therefore remain valid after multiplication by a bounded Borel scalar,
as in \Cref{thm:actual-kernel}.
\end{proof}

\begin{proposition}[Smooth-background energy]
\label[proposition]{prop:first-order-energy}
Let $X/F$ be smooth projective, $F\in\mathcal G_K$, and $L$ ample.
Let $r$ be a bounded classical metric and let $p,q$ be continuous
model-semipositive metrics on the same line. Set $g=r-p$, $b=q-p$
and $V=\deg_F(L^n)$. If $\Omega$ is a product of $n-1$ smooth
semipositive regularized-FS curvatures on $L$, put
\begin{align}
 \Gamma(g)&=\tfrac12D[g^2]-gD[g],\label{eq:gamma}\\
 \Gamma(g,h)&=\tfrac12\bigl(D[gh]-gD[h]-hD[g]\bigr),
 &E_\Omega(g,h)&=\int_{X^{\an}}\Gamma(g,h)\wedge\Omega.\notag
\end{align}
Then
\begin{align}
 E_\Omega(g)&\geq0,&
 |E_\Omega(g,b)|^2&\leq E_\Omega(g)E_\Omega(b),\label{eq:energy-cs}\\
 E_\Omega(g)&=-\int gD[g]\wedge\Omega,&
 E_\Omega(g,b)&=-\int gD[b]\wedge\Omega,\label{eq:energy-green}\\
 0\leq E_\Omega(b)&\leq2V\|q-p\|_\infty.&&\label{eq:energy-bound}
\end{align}
\end{proposition}
\begin{proof}
We first check that the energy expressions are defined using the
first-order measures already constructed. Locally translate the
bounded classical summands to be nonnegative. Their squares are then
classical by convex increasing composition on curves. Polarization
expresses squares and products of their differences as finite linear
combinations of such squares. Thus \Cref{prop:first-order} defines
the differentiated terms in \eqref{eq:gamma}, and its coefficient
measures permit multiplication by the bounded gaps.

The curve chain rule of \Cref{lem:bounded-curve-energy}, transported
through \Cref{thm:actual-kernel}, makes the quadratic expression
nonnegative on each positive decomposable background. A smooth
regularized-FS curvature has a positive semidefinite ambient
Hessian, and their product is strongly positive. With one dominating
coefficient measure, positivity on countably many rational directions
passes by continuity to its variable decomposable directions.
Polarization proves \eqref{eq:energy-cs}.

Stokes against the smooth closed $\Omega$ gives the first equality
of \eqref{eq:energy-green}. For the second, approximate $p,q$
uniformly by smooth regularized-FS metrics, keeping $\Omega$ fixed.
The localized variation estimate of \Cref{lem:curve-variation}
supplies an integrable factor multiplying the uniform approximation
error on the curve slices. A finite partition of unity and the kernel formula
therefore give Green symmetry with the continuous $b$. The positive
continuous measures $c_1(L,p)\wedge\Omega$ and
$c_1(L,q)\wedge\Omega$ each have mass $V$. Their difference gives
\eqref{eq:energy-bound}.
\end{proof}

\begin{remark}[The energy current]\label[remark]{rem:gamma-not-closed}
The current $\Gamma(g)$ need not be closed. The face-tangency
argument above is applied only to $D[v]$. Energy positivity is used
on the strongly positive smooth backgrounds of
\Cref{prop:first-order-energy}. Its construction uses only first-order
currents, independently of the higher mixed products.
\end{remark}

\section{Contact rigidity and decreasing regularization}
\label{sec:contact}

In this section \(X/F\) is smooth projective geometrically integral,
\(F\in\mathcal G_K\), and \(L\) is ample. The continuous products
used here are constructed from model-semipositive metrics. Uniform
pure-FS approximation and smoothing are valid over these fields
\cite[Corollary~7.6 and Theorems~7.8, 7.14]{BE}. Their model measure is
\begin{equation}\label{eq:native-model-measure}
 c_1(L_1,\phi_1)\cdots c_1(L_n,\phi_n)
 =\sum_B b_B\deg_B(\mathcal L_1\cdots\mathcal L_n)\delta_{x_B}
\end{equation}
on a common normal model, where \(b_B\) is the multiplicity of the
special-fiber component \(B\) and all intersection degrees are taken
over the residue field. This convention retains residual degrees
and the prescribed absolute value.

\begin{definition}\label[definition]{def:envelope}
For a continuous metric \(\rho\) on \(L\), let
\[
 P_L\rho=\sup\{s:s\text{ continuous model-semipositive on }L,\ s\leq\rho\},
\]
pointwise in a common local frame. Nef rational model minorants
give the same envelope: approximate a continuous competitor
uniformly and shift its approximants down by rational valuation
constants. Continuity and orthogonality of this envelope are proved
in \Cref{thm:normal-envelope}.
\end{definition}

The next lemma replaces one Fubini--Study curvature by integration
over a generic hyperplane. It is used as an identity of measures:
we will integrate the bounded, possibly discontinuous gap $r-p$
against it in the proof of contact rigidity.

\begin{lemma}[A single generic hyperplane]\label[lemma]{lem:single-hyperplane}
Suppose \(n\geq2\). Let \(q\) be a pure-FS metric on \(L\), given by
generating sections \(s_0,\ldots,s_N\) of \(mL\), with rational
constants absorbed in a tensor power. Put
\[
 E=\widehat{F(Z_0,\ldots,Z_N)}_{\rm Gauss},\qquad
 H=\left\{\sum_iZ_i s_i=0\right\}\subset X_E.
\]
Then \(H/E\) is smooth projective geometrically integral of
dimension \(n-1\). For the natural map
\(\iota:H^{\an}\to X^{\an}\) and continuous model-semipositive
metrics \(a_1,\ldots,a_{n-1}\) on \(L\),
\begin{align}
 c_1(q)\prod_{i=1}^{n-1}c_1(a_i)
 &=\frac1m\,\iota_*\prod_{i=1}^{n-1}c_1(a_i|_H),
 \label{eq:single-hyperplane}\\
 \deg_E((L_E|_H)^{n-1})&=m\deg_F(L^n).
 \label{eq:hyperplane-degree}
\end{align}
The first equality is an equality of Radon measures.
\end{lemma}
\begin{proof}
The associated morphism \(X\to\mathbb P_F^N\) is finite onto its
image: the pullback of \(\mathcal O(1)\) is ample, so a
positive-dimensional proper fiber is impossible. Generic Bertini
in characteristic zero gives smoothness and geometric
irreducibility of the hyperplane pullback; for irreducibility see
\cite[Lemma~37.32.3]{Stacks}. Completion of the Gauss parameter
field preserves these geometric properties.

First take all \(a_i\) to be nef rational model metrics,
and choose a common normal projective model carrying them and the
generated model line of \(q\). Work over
\(R[Z_0,\ldots,Z_N]_{(\pi)}\) and then its completion. For an
irreducible special component \(B\), the incidence in affine
coefficient space is the kernel vector bundle of the generating
evaluation map. Its projectivization is a projective bundle over
\(B\). If the image of \(B\) under the
generated linear system is a point, the generic hyperplane misses
it, and its intersection degree with that generated line is zero.
Otherwise the incidence dominates the parameter space; its generic
fiber is integral. It need not be geometrically integral.

The multiplicity of the induced component is the original
vertical multiplicity \(b_B\). This can be read at its generic
point, away from the singular loci of the model and of \(B\):
the hyperplane cuts a parameter on \(B\), whereas the uniformizer
has order \(b_B\). Such generic points avoid every fixed proper
closed subset of \(B\) in the universal incidence. Normalization
retains the associated residual degrees. The divisorial points of
the hyperplane model restrict to \(x_B\), and the projection formula
for their intersection degrees gives exactly the coefficient of
\(\delta_{x_B}\) in \eqref{eq:native-model-measure}. The root
presentation of \(mL\) supplies the factor \(1/m\).
Completion preserves normality because it is a regular morphism.

This proves \eqref{eq:single-hyperplane} for model metrics.
Uniform approximation passes it to the stated continuous
metrics, on both sides as Radon measures. Finally the algebraic
intersection number of \(H\sim mL_E\) gives
\eqref{eq:hyperplane-degree}. There is no finite-degree factor for
the transcendental Gauss extension.
\end{proof}

\begin{lemma}[Weighted uniform stability in lower dimension]
\label[lemma]{lem:weighted-stability}
Suppose that bounded classical metrics in dimension \(d\) admit
decreasing continuous semipositive regularization. On a smooth
projective \(d\)-fold, let \(r\) be a bounded classical metric,
\(p\) continuous model-semipositive on the same ample line, and
\(h=r-p\). For continuous semipositive \(a_i,b_i\) on that line,
with \(W\) its top degree,
\begin{equation}\label{eq:weighted-stability}
 \left|\int h\left(\prod_{i=1}^d c_1(a_i)
                         -\prod_{i=1}^d c_1(b_i)\right)\right|
 \leq2W\sum_{i=1}^d\|a_i-b_i\|_\infty.
\end{equation}
\end{lemma}
\begin{proof}
Choose continuous semipositive \(r_\ell\downarrow r\).
At a fixed \(\ell\), telescope the product and use the continuous
Green formula to move the difference
\(D[a_i-b_i]\) onto \(r_\ell-p\). The two resulting positive
mixed measures each have mass \(W\). The corresponding term is
bounded by \(2W\|a_i-b_i\|_\infty\). Sum the terms and let
\(\ell\to\infty\). On the two fixed measures in the left side,
the functions \(r_\ell-p\) are uniformly bounded and converge
pointwise to \(h\), so bounded convergence applies.
\end{proof}

\begin{theorem}[Projective contact and regularization]
\label[theorem]{thm:projective-regularization}
For every \(F\in\mathcal G_K\) and every positive dimension, the
following statements hold.
\begin{enumerate}
\item If \(r\) is a bounded classical metric on \(L\), \(p\) is
continuous model-semipositive, \(r\geq p\), and
\[
 \int (r-p)c_1(L,p)^n=0,
\]
then \(r=p\) at every point.
\item Every bounded classical metric \(r\) on \(L\) is the
decreasing pointwise limit of continuous model-semipositive
metrics, and also of rational pure-FS metrics, on the same
identified line over \(F\).
\end{enumerate}
\end{theorem}
\begin{proof}
We prove contact in dimension $n$ using both assertions in smaller
dimensions, and then deduce regularization in dimension $n$ from
contact. The induction is simultaneous for all $F\in\mathcal G_K$,
so the lower-dimensional assertions remain available over the
unit-Gauss field of a generic hyperplane. In dimension one, contact
is \Cref{prop:curve-contact}; the envelope argument at the end of the
proof then gives regularization.

Assume \(n\geq2\), and set \(g=r-p\geq0\), \(V=\deg_F(L^n)\).
Choose smooth semipositive regularized-FS metrics \(p_j\to p\)
uniformly and put \(\Omega_j=c_1(p_j)^{n-1}\). The function \(g\)
is bounded and upper semicontinuous. Continuous mixed convergence
and the upper Portmanteau inequality give
\[
 A_j=\int g\,c_1(p)\wedge\Omega_j\longrightarrow0.
\]
Since $D[g]=c_1(r)-c_1(p)$, the Green identity gives
$E_{\Omega_j}(g)=A_j-\int g\,c_1(r)\wedge\Omega_j$.
The second term is nonnegative by first-order positivity and $g\geq0$.
Together with \Cref{prop:first-order-energy}, this yields
\[
 0\leq E_{\Omega_j}(g)\leq A_j.
\]
For a fixed pure-FS metric \(q\), set
\(I_j(q)=\int g\,c_1(q)\wedge\Omega_j\).
Green symmetry identifies $I_j(q)-A_j$ with
$-E_{\Omega_j}(g,q-p)$. Cauchy--Schwarz and
$E_{\Omega_j}(q-p)\leq2V\|q-p\|_\infty$ therefore yield
\begin{equation}\label{eq:contact-estimate}
 0\leq I_j(q)
 \leq A_j+\bigl(2V\|q-p\|_\infty A_j\bigr)^{1/2}
 \longrightarrow0.
\end{equation}

The convergence of the background measures alone does not identify
their integrals against the discontinuous function $g$. Apply
\Cref{lem:single-hyperplane} to $q$. On its hyperplane $H/E$, the
dimension-$(n-1)$ induction gives decreasing regularization of $r|_H$,
which is the hypothesis needed for the weighted stability estimate. Applying
\Cref{lem:weighted-stability} there and using
\eqref{eq:hyperplane-degree}, the factor \(m\) cancels the \(1/m\)
in \eqref{eq:single-hyperplane}. Consequently
\[
 I_j(q)\longrightarrow
       \int g\,c_1(q)c_1(p)^{n-1}.
\]
The measure identity and
\eqref{eq:contact-estimate} now give
\[
 \int_{H^{\an}}g|_H\,c_1(p|_H)^{n-1}=0.
\]
The lower-dimensional contact theorem shows \(g|_H=0\).

We now pass from vanishing on the generic hyperplane to pointwise
vanishing on $X$. The intermediate points are the Shilov boundary
points of a model associated with a pure-FS metric $q$ close to $p$.
Normalize the projective closure of its image in the finite
function-field extension \(F(X)\). This gives a normal model
finite over that image, with the generated model line of \(mL\)
relatively ample. Every special component has dimension \(n\)
and meets the generic hyperplane. The divisorial points of these
intersections restrict to all its component points \(x_B\), so
the equality \(g|_H=0\) implies \(g(x_B)=0\).

On a formal generating-frame domain, the model weight of \(mq\)
is zero. The local weight \(m(r-q)\) is classical psh there.
The maximum principle on its Shilov boundary
\cite[Theorem~6.4 and Remark~7.5]{GR} therefore gives
\[
 r-q\leq\max_B(p-q)(x_B)\leq\|p-q\|_\infty.
\]
These domains cover \(X^{\an}\). Hence
\[
 0\leq r-p\leq2\|p-q\|_\infty.
\]
Pure-FS approximation to \(p\) proves the contact assertion.

It remains to deduce regularization from contact. Compact
metrizability gives continuous metrics $\rho_\ell\downarrow r$.
Set $s_\ell=P_L\rho_\ell$. By \Cref{thm:normal-envelope}, each
$s_\ell$ is continuous model-semipositive and satisfies envelope
orthogonality. Although $r$ is not yet known to be a semipositive
competitor, the classical metric $\max(r,s_\ell)$ lies above $s_\ell$
and satisfies
\[
 0\leq\max(r,s_\ell)-s_\ell\leq\rho_\ell-s_\ell.
\]
Its contact integral is therefore zero. The contact assertion just
proved implies $\max(r,s_\ell)=s_\ell$, or $r\leq s_\ell$.
Since \(s_\ell\leq\rho_\ell\) and envelopes preserve order,
\[
 s_\ell\downarrow r.
\]
Uniform pure-FS approximation need not preserve the order of the
sequence. To retain monotonicity, choose rational pure-FS $a_\ell$ with
\(\|a_\ell-s_\ell\|_\infty\leq\epsilon_\ell\), where
\(\epsilon_\ell=\lambda4^{-\ell}\).
Then \(q_\ell=a_\ell+3\epsilon_\ell\) is decreasing: its next
term is at most \(s_\ell+\epsilon_\ell\), whereas its current
term is at least \(s_\ell+2\epsilon_\ell\).
It converges pointwise to \(r\). This proves the second assertion
and completes the simultaneous induction.
\end{proof}

\begin{proposition}[Local algebraic regularization]
\label[proposition]{prop:local-algebraic-regularization}
Let \(F\in\mathcal G_K\), and let \(u\) be locally bounded
classical psh on an ordinary open of a smooth geometrically integral
algebraic \(F\)-variety. Near each point, \(u\) is the decreasing limit
of rational semipositive PL functions over \(F\) on one fixed
ordinary domain. The approximants may be taken in Wu's algebraic
tropically convex PL class on that domain.
\end{proposition}
\begin{proof}
Choose an affine algebraic neighborhood of the prescribed point.
Compactify it projectively over \(F\), and resolve the compactification
while preserving the smooth affine open
\cite[Theorem~1.1]{BM}. The resulting variety is smooth projective
and geometrically integral over \(F\), and contains the ordinary open
on which \(u\) is given.
The projective theorem applies to a global metric, whereas $u$ is
given only locally. Apply \Cref{lem:common-collar} to the single
function $u$. On a smaller fixed neighborhood it realizes $u+c$
as the exact weight of a bounded classical metric on an ample line;
outside a larger neighborhood the metric is supplied by a fixed
Fubini--Study barrier.
\Cref{thm:projective-regularization} provides decreasing
pure-FS metrics there. Restrict to the fixed frame and subtract
the same constant \(c\). The inactive floor of
\Cref{lem:inactive-floor} gives the stated finite algebraic
tropically convex representations, including at zeros of
individual sections. 
\end{proof}

\section{Coefficient measures and the mixed-current construction}
\label{sec:construction}

We construct the products by induction using
$T_k=D(u_1T_{k-1})$. Two points require attention: the Borel product
$u_1T_{k-1}$ must be independent of a moment presentation, and
top-degree PL convergence must determine the pairing with an arbitrary
complementary test. The coefficient-measure lemma addresses the first
point and supplies the linear-algebra reduction for the second.

\begin{lemma}[Coefficient measures of strongly positive currents]
\label[lemma]{lem:strong-coefficients}
Every symmetric strongly positive original current is of
coefficient Radon type. On a fixed finite moment presentation its
coefficient variations are bounded by finitely many positive
cutoff evaluations. Consequently these variations are locally
uniformly bounded in a weakly convergent family of such currents.
\end{lemma}
\begin{proof}
Put \(r=n-k\), and let \(E\) be the ambient moment tangent space.
Identify symmetric \((r,r)\)-forms with symmetric bilinear tensors
using the usual Lagerberg sign. Choose a rational basis \(e_I\)
of \(\bigwedge^rE^*\), and put
\(\sigma=\sum_Ie_I\otimes e_I\). For any constant symmetric
tensor \(A\), choose \(C_A\) so that \(C_A\sigma\pm A\) are
positive semidefinite, hence weakly positive. If \(\phi\) is an
arbitrary original smooth scalar supported where \(\kappa=1\),
strong positivity gives
\begin{equation}\label{eq:coefficient-bound}
 |\langle T,\phi A\rangle|
 \leq C_A\|\phi\|_\infty\langle T,\kappa\sigma\rangle.
\end{equation}
The estimate applies to every original smooth scalar $\phi$,
including those that do not factor through the chosen moment map:
pass to a common refinement, where pullback preserves weak positivity.
Riesz representation therefore gives coefficient measures on the
analytic domain itself.

Let \(W_r\) be the span of squares of decomposable \(r\)-vectors
in \(\operatorname{Sym}^2(\bigwedge^rE)\), and let \(A_r=W_r^\perp\)
in the dual symmetric tensor space. Both signs of a tensor in
\(A_r\) are weakly positive. Hence \(T\) annihilates every smooth
coefficient times such a tensor. Its coefficient measure is zero,
and remains zero after multiplication by any bounded Borel
function.

We use decomposable directions after passing to the quotient by
$A_r$. Choose finitely many rational decomposable squares and a fixed
linear splitting in that quotient. Every symmetric test then has a
decomposition
\begin{equation}\label{eq:full-test-splitting}
 \eta=\sum_\ell a_\ell\Omega_\ell+B,\qquad B(t)\in A_r,
\end{equation}
where the \(a_\ell\) are original smooth compact scalars and the
\(\Omega_\ell\) are constant rational decomposable forms.
The annihilator term is killed by \(T\) and by every Borel
multiple of \(T\).

All form relations and changes of presentation can be checked
after multiplying by every original smooth scalar. Uniqueness
of Radon measures then proves those relations as measure
identities. This also proves compatibility with Borel
multiplication. Finally, in a weakly convergent family, the
finitely many positive evaluations on \(\kappa\sigma\) converge
and are bounded. Equation~\eqref{eq:coefficient-bound} gives the
last assertion.
\end{proof}

\begin{remark}[The middle-degree annihilator]
\label[remark]{rem:pluecker-annihilator}
The term \(B\) in \eqref{eq:full-test-splitting} is necessary.
For \(\dim E=4\) and \(r=2\), the Plücker quadratic
\[
 z_{12}z_{34}-z_{13}z_{24}+z_{14}z_{23}
\]
vanishes on all decomposable bivectors. The decomposable-square
span has dimension \(20\), whereas the symmetric tensor space has
dimension \(21\). Strong positivity kills this annihilator in
the current pairing; it does not remove it from the original
form sheaf.
\end{remark}

To apply a top-degree convergence theorem to a test in complementary
degree $r$, we will fill its $r$ directions with curvature factors.
The following quadratic potentials realize exactly those factors and
admit PL approximants on the same domain.

\begin{lemma}[Auxiliary quadratic potentials]\label[lemma]{lem:quadratics}
After fixed finite localization, every decomposable direction in
\eqref{eq:full-test-splitting} is a product
\[
 \Omega=\prod_{a=1}^{r}D h_a,\qquad
 h_a=\tfrac12\ell_a^2,
\]
where the \(\ell_a\) are rational linear combinations of
logarithmic norms of algebraic units. Each \(h_a\) admits a
decreasing sequence of algebraic tropically convex rational PL
approximants, uniformly on the relevant compact neighborhood.
\end{lemma}
\begin{proof}
Use \Cref{lem:exact-units} and clear rational denominators.
The unit logarithms are pluriharmonic, so
\(D(\ell_a^2/2)=d'\ell_a\wedge d''\ell_a\).
On one fixed bounded interval containing the range of \(\ell_a\),
approximate \(t^2/2\) uniformly by finite maxima of affine
functions with rational slopes and constants in
\(\lambda\mathbb Q\). Such maxima are dense: approximate finitely
many supporting slopes and their intercepts, using the density
of \(\mathbb Q\) and of \(\lambda\mathbb Q\).
Choose errors \(\epsilon_m=\lambda4^{-m}\) and add
\(3\epsilon_m\). The same margin calculation as in
\Cref{thm:projective-regularization} makes the sequence decreasing
to \(t^2/2\). Clearing denominators realizes the affine branches
by algebraic monomials in the units and powers of \(\pi\).
\Cref{lem:inactive-floor} gives the required algebraic PL
presentation. The interval, units and analytic domain remain fixed.
\end{proof}

\begin{proposition}[Top-degree PL convergence and normalization]
\label[proposition]{prop:wu-input}
On an ordinary algebraic analytic open of dimension \(n\) over
\(K\), decreasing sequences of algebraic tropically convex PL functions
\(p_{i,j}\downarrow p_i\), \(0\leq i\leq n\), with locally
bounded limits satisfy
\begin{equation}\label{eq:wu-weighted}
 p_{0,j}\,\MA_{\rm PL}(p_{1,j},\ldots,p_{n,j})
 \longrightarrow
 p_0\,\MA_{\rm PL}(p_1,\ldots,p_n)
\end{equation}
locally weakly as Radon measures. The continuous PL
stages agree with the original CLD products normalized by
\eqref{eq:native-model-measure}.
\end{proposition}
\begin{proof}
This is \cite[Theorem~4.2(4)]{Wu}, applied on compact supports,
with every sequence, including the undifferentiated one, in the
algebraic tropically convex PL class.
The local class is \cite[Definition~3.8]{Wu}; the continuous PL
comparison is \cite[Proposition~4.1(3)]{Wu}.

We make the normalization explicit. Wu's preliminary convention
may be achieved by replacing the absolute value temporarily by
\(|\cdot|_*=|\cdot|^{1/\lambda}\). The canonical homeomorphism
\(j:X^{\an}\to X_*^{\an}\) sends a seminorm to its
\(1/\lambda\)-power. Put \(p_{i,*}\circ j=p_i/\lambda\).
In logarithmic coordinates \(t_*=t/\lambda\), pullback of a
top Lagerberg form satisfies
\begin{equation}\label{eq:valuation-calibration}
 \int_{X^{\an}}j^*\alpha
       =\lambda^{-n}\int_{X_*^{\an}}\alpha.
\end{equation}
There are two differentials in each of the \(n\) directions,
while the calibrated real volume has dimension \(n\), which
gives this factor. Also
\(j^*D_*p_{i,*}=\lambda^{-1}Dp_i\).
The factors therefore cancel in the calibrated top product:
\[
 j^*\MA_*(p_{1,*},\ldots,p_{n,*})
       =\MA_{\rm PL}(p_1,\ldots,p_n)
\]
as transported measures. Equivalently, both sides assign the
same structural degree to every model component in
\eqref{eq:native-model-measure}. Inserting the weight gives
the factor \(\lambda^{-1}\), since \(p_{0,*}\circ j=p_0/\lambda\);
multiplication by \(\lambda\) returns \eqref{eq:wu-weighted}.
Thus the formula holds for the prescribed absolute value and the given
functions, with the model masses in \eqref{eq:native-model-measure}.

The constructed PL functions satisfy the stated local hypothesis
by \Cref{lem:inactive-floor,lem:quadratics}. In particular rational
valuation constants are encoded by constant algebraic
coordinates, or become rational under the normalization above.
All uses of \eqref{eq:wu-weighted} are on compact tests; no global
integrability assertion on a noncompact open is needed.
\end{proof}

\begin{proposition}[Construction and decreasing PL convergence]
\label[proposition]{prop:pl-construction}
On an ordinary open in the analytification of a smooth geometrically
integral \(K\)-variety, the
currents in \Cref{thm:mixed-currents} exist. If each potential varies
in a decreasing sequence of algebraic tropically convex PL functions
on one fixed open, with locally bounded classical limit, then the
weighted currents converge on every original complementary test.
\end{proposition}
\begin{proof}
We prove existence in degree $k$ from weighted convergence in degree
$k-1$, and then prove weighted PL convergence in degree $k$.
This order is required by the recursion: differentiation of
$u_1T_{k-1}$ uses the undifferentiated argument at the preceding degree.
For $k=0$, $T_0=[1]$, and the weighted assertion is bounded convergence
against each fixed original top-form measure.

Suppose the statement holds in degree \(k-1\). The coefficient
measures of \(T_{k-1}\) allow multiplication by \(u_1\).
Distributional differentiation of this Borel product defines \(T_k\) by
\eqref{eq:recursive-current}. Choose decreasing rational
pure-FS approximants to its potentials on one fixed open.
By weighted convergence in degree $k-1$, the Borel products formed
from these approximants converge to $u_1T_{k-1}$ on every original
test. Applying distributional differentiation shows that their
continuous mixed currents converge to the $T_k$ just defined. Finite convex-maximum smoothing makes
each approximating current strongly positive. The limit is
strongly positive, and \Cref{lem:strong-coefficients} supplies
its coefficient measures. Permuting the continuous
approximants proves symmetry. Closedness and multilinearity
follow from the recursion and the continuous identities.
The recursion also makes the candidate independent of the
chosen approximants.

We still need weighted convergence in degree $k$; it does not follow
by multiplying the unweighted current limit by a discontinuous
function. Use \eqref{eq:full-test-splitting}. Its annihilator term
has zero coefficient measure, so both the current and its Borel
multiple kill that term. It suffices to treat a test \(a\Omega\), where
\(\Omega=Dh_1\cdots Dh_r\), \(r=n-k\), and the \(h_\alpha\)
are as in \Cref{lem:quadratics}. First we prove the unweighted
identity of Radon measures on the analytic domain
\begin{equation}\label{eq:current-pl-measure}
 T_k(u_1,\ldots,u_k)\wedge\Omega
 =\MA_{\rm PL}(u_1,\ldots,u_k,h_1,\ldots,h_r).
\end{equation}
Let \(u_{i,j}\downarrow u_i\) be PL sequences and
\(h_{\alpha,l}\downarrow h_\alpha\) their uniform PL
approximants. At fixed \(j\), continuous mixed convergence lets
one choose an auxiliary index \(l(j)\geq j\), increasing in
\(j\), so that the error on the given smooth scalar test is
less than \(1/j\). Apply \Cref{prop:wu-input} to the resulting
\(n\) decreasing PL sequences. Its unweighted version
follows by taking a fixed nonzero weight in \(\lambda\mathbb Q\).
The other limit is
\(T_k(u_1,\ldots,u_k)\wedge\Omega\), by the already established
current convergence. Equality on every smooth scalar, and
Radon uniqueness, prove \eqref{eq:current-pl-measure}.

Because \eqref{eq:current-pl-measure} is now an identity of Radon
measures on the analytic domain, we may multiply it by the given
Borel function $u_0$. For the varying PL weight $u_{0,j}$, at
each fixed \(j\) use the continuous compact scalar
\(a u_{0,j}\) to choose the auxiliary index with error less
than \(1/j\). Increase the indices to preserve all
decreasing auxiliary sequences. The weighted assertion
\eqref{eq:wu-weighted} yields
\[
 \langle u_{0,j}T_k(u_{1,j},\ldots,u_{k,j}),a\Omega\rangle
 \longrightarrow
 \int a u_0\,d\MA_{\rm PL}(u_1,\ldots,u_k,h_1,\ldots,h_r).
\]
By \eqref{eq:current-pl-measure}, the right side is precisely
\(\langle u_0T_k(u_1,\ldots,u_k),a\Omega\rangle\).
For each test, this diagonal proves convergence to the current
defined by the recursion. Finite summation finishes
the induction.
\end{proof}

\section{Decreasing families of classical potentials}\label{sec:arbitrary-columns}

Each classical potential has a decreasing PL regularization, but
independently regularizing the members of a decreasing family need
not preserve their order. We use a compact diagonal argument to
choose one approximant from each row. The choice must both preserve
decrease and satisfy prescribed bounds on current pairings.

\begin{lemma}[A compact triangular diagonal]
\label[lemma]{lem:triangular-diagonal}
Let \(Z\) be compact, and suppose, for finitely many \(i\), that
bounded metrics \(R_{i,j}\downarrow R_i\) on one identified line
have continuous decreasing rows \(P_{i,j,m}\downarrow R_{i,j}\).
Let \(\delta_j>0\) decrease strictly to zero. Given arbitrary
thresholds for \(m\), one can choose \(m_j\geq j\) beyond them
so that
\[
 Q_{i,j}=P_{i,j,m_j}+\delta_j
\]
decreases in \(j\) and converges pointwise to \(R_i\).
\end{lemma}
\begin{proof}
Proceed inductively. The desired inequality
\(P_{i,j,m}+\delta_j<Q_{i,j-1}\) holds eventually at each
point: its limit on the left is
\(R_{i,j}+\delta_j\), whereas the right is at least
\(R_{i,j-1}+\delta_{j-1}\), with a strict positive margin.
The sets on which it holds are increasing open sets in \(m\).
Compactness gives one index for all points and the finite
set of indices \(i\).

Monotonicity alone would only give a limit above $R_i$. To force the
correct limit, compare the chosen term with the first $j$ fixed rows
by imposing the finite family of caps
\[
 P_{i,j,m}+\delta_j\leq P_{i,a,j}+2\delta_j,
 \qquad 1\leq a\leq j.
\]
The limiting left side is at most
\(R_{i,a}+\delta_j\), so the same strict-margin compactness
argument permits all these caps. Include the prescribed
thresholds and \(m_j\geq j\).
The decreasing limit \(Q_i\) is at least \(R_i\).
For fixed \(a\), the cap gives \(Q_i\leq R_{i,a}\) as
\(j\to\infty\). Letting \(a\to\infty\) gives \(Q_i=R_i\)
pointwise.
\end{proof}

\begin{proposition}[Convergence for decreasing classical potentials]
\label[proposition]{prop:arbitrary-convergence}
The conclusion of \Cref{thm:joint-monotone} holds on ordinary
opens in the analytification of a smooth geometrically integral
\(K\)-variety.
\end{proposition}
\begin{proof}
The decreasing-limit property on curves shows that each locally
bounded limit $u_i$ is classical
\cite[Proposition~4.2]{GR}.
The endpoint inequalities
\(u_i\leq u_{i,j}\leq u_{i,1}\) give common upper and lower
bounds near a prescribed compact local support. Apply
\Cref{lem:common-collar} on a fixed smooth projective compactification
of the affine algebraic chart, as in the proof of
\Cref{prop:local-algebraic-regularization}, before choosing any regularizing
sequences. It gives one smooth projective variety \(Z\), one
ample line, one frame on a relatively compact ordinary
neighborhood \(V\), and one constant \(c\), with bounded
classical metrics
\[
 R_{i,j}\downarrow R_i
\]
whose local weights are exactly \(u_{i,j}+c\) and \(u_i+c\).
For each \(i,j\), choose decreasing pure-FS metrics
\(P_{i,j,m}\downarrow R_{i,j}\)
by \Cref{thm:projective-regularization}. One inactive floor
works for all these sequences on \(V\).

Fix a complementary test \(\eta\) supported in \(V\). Write
\[
 F_j=\langle S_k(u_{0,j};u_{1,j},\ldots,u_{k,j}),\eta\rangle,
 \qquad
 B_j=|\langle T_k(u_{1,j},\ldots,u_{k,j}),\eta\rangle|.
 \]
Choose strictly decreasing \(\delta_j\in\lambda\mathbb Q_{>0}\)
with \(\delta_j\leq2^{-j}\) and \(\delta_jB_j\leq1/j\).
The small shift enforces strict decrease in the compact diagonal.
It also changes the undifferentiated argument, so we choose it to make
that change in the pairing at most $1/j$. This choice uses only the
finite value of $B_j$, without a bound uniform in $j$.
For each fixed \(j\), decreasing PL convergence from
\Cref{prop:pl-construction} shows that, as \(m\to\infty\),
the weighted-current pairing of the shifted local weights
\(P_{i,j,m}+\delta_j-c\) converges to
\[
 F_j+\delta_j\langle T_k(u_{1,j},\ldots,u_{k,j}),\eta\rangle.
\]
The shifts in the differentiated arguments have zero curvature.
Choose a threshold for \(m\) making this error less than \(1/j\).

Apply \Cref{lem:triangular-diagonal} with those thresholds.
The caps in the diagonal lemma ensure that the local weights
$Q_{i,j}-c$ converge pointwise to $u_i$. They are decreasing
algebraic tropically convex PL functions on the same $V$. Their weighted-current pairing \(G_j\) converges
to the desired limit by \Cref{prop:pl-construction}.
The chosen threshold and the bound on \(\delta_j\) give
\[
 |F_j-G_j|\leq2/j.
\]
This proves convergence on \(\eta\). A finite original
smooth partition of the support proves the assertion for
an arbitrary compact test. The unweighted assertion follows
by taking a constant weight.
\end{proof}

\begin{proposition}[Continuous compatibility and top degree]
\label[proposition]{prop:continuous-top}
The constructed currents agree with the continuous CLD
products on their common continuous domain. Their top-degree
weighted convergence holds against continuous compact scalar
tests.
\end{proposition}
\begin{proof}
For continuous limits, decreasing continuous PL
regularization is uniform on compact subsets by Dini's
argument. Continuous CLD mixed convergence therefore
identifies the products. Conversely the recursive formula
and \Cref{prop:pl-construction} determine the current
uniquely from any decreasing PL regularization.

In top degree choose a nonnegative original smooth cutoff
equal to one on the test support. Its evaluations on the
positive measures \(T_{n,j}\) converge by
\Cref{prop:arbitrary-convergence}, giving a common mass
bound there. The weights have a common bound \(M\)
on a slightly larger compact neighborhood, so
\(|S_{n,j}|\leq M T_{n,j}\) there. Uniform approximation
of a continuous compact scalar by original smooth scalars
then extends convergence to that scalar.
\end{proof}

\section{Ordinary localization and the main theorems}
\label{sec:localization}

The preceding constructions were made on algebraic charts. We now
transfer them to the analytic spaces of \Cref{conv:spaces}. The
ordinary neighborhood must be fixed before regularization, and a
finite change in the constant field must retain the original
integration calibration. Once these points are settled, the recursive
formula identifies the locally constructed currents on overlaps.

\begin{proposition}[Transport to an ordinary analytic chart]
\label[proposition]{prop:ordinary-transport}
Every point of \(X\) has one ordinary neighborhood on which
\Cref{prop:local-algebraic-regularization,prop:pl-construction,prop:arbitrary-convergence}
apply with the original \(K\)-calibration. For a finite decreasing
family of locally bounded classical functions, the neighborhood
and the projective variety can be chosen independently of the
sequence index.
\end{proposition}
\begin{proof}
By \Cref{lem:local-algebraization}, a sufficiently small ordinary
neighborhood is an ordinary open of a smooth affine \(K\)-variety.
Its integral components are disjoint; restrict to the component
containing the point. Let \(L\) be its finite constant field as
in \Cref{lem:finite-constants}. The same variety is smooth
geometrically integral over the finite \(p\)-adic field \(L\).
Take its projective closure and resolve that closure while preserving
the smooth affine open \cite[Theorem~1.1]{BM}. This gives a smooth
projective geometrically integral \(L\)-variety containing the
original chart. 
The preceding theorems apply with \(L\) as the initial field and
\(\mathcal G_L\) as its auxiliary tower class. All local weights
are functions on the identical original analytic chart.
The algebraic \(L\)-functions used in their PL expressions belong
to its original \(K\)-coordinate ring. If \(\pi_L\) is a uniformizer
of \(L\), then
\(-\log|\pi_L|\mathbb Q=\lambda\mathbb Q\).
Thus these are rational PL functions defined over \(K\), and the
classical and model-semipositive conditions are preserved.
For the currents, multiply the \(L\)-calibration by the
single factor \([L:K]\) in
\eqref{eq:finite-constant-calibration}.

For a decreasing family \(u_{i,j}\downarrow u_i\), choose this
chart and a slightly larger compact neighborhood first.
The endpoint inequalities
\[
 u_i\leq u_{i,j}\leq u_{i,1}
\]
give common two-sided bounds there for the finite family.
The common collar of \Cref{lem:common-collar} therefore fixes
the projective variety, line, frame and local domain for all indices.
The potentials themselves are not required to extend to the
compactification. The analytic isomorphism transfers every
original test presentation, so it also preserves the current identities on every
original test form.
\end{proof}

\phantomsection\label{proof:local-regularization}
\begin{proof}[Proof of \Cref{thm:local-regularization}]
The inclusion \(\Cb(U)\subset\Rb(U)\) follows from
\Cref{prop:ordinary-transport}.
Conversely, semipositive PL functions are classical by
\cite[Theorem~4.3]{GR}. This result is applied on the strict
ordinary algebraic charts just constructed; rational tensor
denominators can be cleared.
The decreasing-limit property of classical subharmonic functions
on curves, with the all-extension quantifier, makes the pointwise
limit classical \cite[Proposition~4.2]{GR}.
That limit is upper semicontinuous as an infimum of continuous
functions, and it is locally bounded by the definition
of \(\Rb\). This proves equality on every ordinary open.

Now let \(u_j\) decrease in either class with a locally
bounded pointwise limit \(u\). The equality just proved makes
every \(u_j\) classical. Since
\(\{u<a\}=\bigcup_j\{u_j<a\}\), the limit is upper semicontinuous.
The curve decreasing-limit property again makes \(u\) classical.
Applying the regularization assertion to this function gives
membership in \(\Rb\) on one neighborhood at a time. No
intersection of the separate approximation neighborhoods
of the \(u_j\) is used.
\end{proof}

\begin{lemma}[Equality and Borel multiplication]
\label[lemma]{lem:borel-equality}
If two coefficient Radon currents \(T,T'\) agree on all
original smooth compact tests on an ordinary open, then
\(hT=hT'\) for every locally bounded Borel function \(h\).
\end{lemma}
\begin{proof}
For a fixed compact original test \(\eta\), the scalar
functionals
\[
 a\longmapsto\langle T,a\eta\rangle,\qquad
 a\longmapsto\langle T',a\eta\rangle
\]
agree on every original smooth scalar. Smooth scalar density
on a compact neighborhood and uniqueness of Radon measures
give \(\mu_{T,\eta}=\mu_{T',\eta}\). Integrating the given
function \(h\) against this equality proves the assertion.
The coefficient relations and common refinements are measure
identities by \Cref{lem:strong-coefficients}; finite localization
therefore removes the choice of presentation.
\end{proof}

\phantomsection\label{proof:main-currents}
\begin{proof}[Proof of \Cref{thm:mixed-currents,thm:joint-monotone}]
Construct the currents on the ordinary charts of
\Cref{prop:ordinary-transport}. At degree zero their
normalization is the intrinsic integration current \([1]_K\),
by \Cref{lem:finite-constants}. Suppose the constructions
agree in degree \(k-1\) on an overlap. By
\Cref{lem:borel-equality}, multiplication by the same Borel
function \(u_1\) agrees there. Applying \(D\) identifies
the two degree-\(k\) currents in
\eqref{eq:recursive-current}. A second application of the
lemma with \(u_0\) identifies the weighted currents.
Thus the local currents glue uniquely.

Closedness, strong positivity, symmetry and additivity in the
differentiated arguments are local properties, proved in
\Cref{prop:pl-construction}. Positive homogeneity follows
from the same recursion and approximation. A constant in a
differentiated argument contributes zero because the preceding
current is closed. Linearity of Borel multiplication
gives the identity for the undifferentiated argument in \Cref{rem:source-slot}.
The recurrence and the continuous compatibility of
\Cref{prop:continuous-top} give independence of PL
regularization.

For joint convergence, \Cref{thm:local-regularization} first
makes each bounded limit \(u_i\) classical. Cover the compact
support of a test by finitely many ordinary charts and split
the test into compact summands by an original smooth
partition of unity. The compact-support partition theorem
\cite[Corollary~15.2.9]{CLD25} is sufficient; no global
paracompactness assumption is needed. On each slightly larger
compact subchart the bounded endpoints give the common bounds
used in \Cref{prop:arbitrary-convergence}. Apply that proposition,
restore the single constant-field factor, and sum.
Only the test forms are partitioned.
Taking a constant weight proves the unweighted convergence.
In top degree, the cutoff mass argument in
\Cref{prop:continuous-top} extends convergence to all continuous
compact scalar tests.
\end{proof}

\begin{proof}[Proof of \Cref{cor:quasiprojective}]
The analytification of a smooth separated variety over \(K\)
is good, Hausdorff, quasi-smooth and without relative analytic
boundary \cite[Exercise~5.1.1(iii), Fact~5.1.5]{Tem}.
These properties are preserved on ordinary opens.
Every compact test uses only finitely many relatively compact
charts inside the given open, to which the main theorems apply.
\end{proof}

\begin{remark}[Two different boundaries]\label[remark]{rem:boundaries}
If \(Y\subset\overline Y\) is a projective compactification,
the deleted algebraic locus \(\overline Y\setminus Y\)
differs from the relative analytic boundary of
\(Y^{\an}\to\mathcal M(K)\), which is empty for smooth \(Y\).
A potential can be unbounded near the deleted algebraic
locus and remain locally bounded at every point of \(Y^{\an}\).
The present theorems apply on that ordinary open and make no
claim about extension of the potential to \(\overline Y^{\an}\).
For a singular variety they apply on its smooth locus.
\end{remark}

\section{Compatibility for specified potential tuples}
\label{sec:etale}

For a specified tuple of potentials, composition with an \'etale
map gives another tuple of locally bounded classical potentials.
The next theorem identifies its mixed current with the limit of the
continuous-stage pullbacks. Thus the operation is defined through
the potentials and their decreasing regularizations; no pullback of
an arbitrary distributional current is required.

\begin{theorem}[\'Etale compatibility]
\label[theorem]{thm:etale-compatibility}
Let \(f:Y\to X\) be étale between spaces satisfying
\Cref{conv:spaces}. Let \(W\subset X\) be an ordinary open
on which a finite tuple \(u_0,\ldots,u_k\in\Cb(W)\) has
decreasing continuous semipositive regularizations
\(a_{i,m}\downarrow u_i\). Then, on every original compact
test in \(f^{-1}(W)\),
\begin{equation}\label{eq:etale-tuple}
 S_{k,Y}(a_{0,m}\circ f;a_{1,m}\circ f,\ldots,a_{k,m}\circ f)
 \longrightarrow
 S_{k,Y}(u_0\circ f;u_1\circ f,\ldots,u_k\circ f).
\end{equation}
This extension of the continuous-stage operation is independent
of the regularizations and commutes with ordinary restriction
and composition of étale maps.
\end{theorem}
\begin{proof}
Composition of curve tests preserves classicality, and
continuity of \(f\) preserves local boundedness.
At a continuous stage, local uniform smooth semipositive
approximation and the pullback identity for smooth forms
identify the pulled-back operation with the continuous mixed
product of the pulled-back weights. To obtain the required smooth
approximations, apply \Cref{thm:local-regularization} to a continuous
weight. Its decreasing PL approximation is uniform on a smaller
compact neighborhood by Dini's theorem. In the pure-FS
presentations furnished by the proof, smoothing the convex maximum
gives original smooth semipositive weights with the same uniform
limit; see \cite[Theorem~7.14]{BE}. The argument is local on both
spaces.
\Cref{thm:joint-monotone} on \(Y\) gives
\eqref{eq:etale-tuple} with the Borel product on its right,
for every choice of base regularization.
For an étale map \(g:Z\to Y\), pull the same continuous
approximants twice and apply the theorem on \(Z\).
Both operations give the weighted current of
\((u_i\circ f\circ g)_i\). Restriction follows from locality.
\end{proof}

\begin{proposition}[The finite-étale degree formula]
\label[proposition]{prop:etale-degree}
Let \(f:Y\to X\) be finite étale of constant degree \(d\).
For every tuple \(u_0,\ldots,u_k\in\Cb(X)\),
\begin{align}
 f_*S_{k,Y}(u_0\circ f;u_1\circ f,\ldots,u_k\circ f)
 &=dS_{k,X}(u_0;u_1,\ldots,u_k),\label{eq:etale-source-degree}\\
 f_*T_{k,Y}(u_1\circ f,\ldots,u_k\circ f)
 &=dT_{k,X}(u_1,\ldots,u_k).\label{eq:etale-curvature-degree}
\end{align}
For locally constant degree, the formulas hold componentwise.
\end{proposition}
\begin{proof}
For smooth inputs, apply the calibrated finite-map integration
formula to the pullback of
\(u_0Du_1\cdots Du_k\wedge\eta\);
see \cite[Proposition~15.4.3]{CLD25}.
Locally uniform smooth approximation gives the same formula
for continuous semipositive inputs.
For the bounded tuple, split a compact base test into finitely
many compact summands supported on ordinary opens \(W_a\)
where simultaneous decreasing continuous regularizations
exist. Properness of \(f\) makes the pulled-back test compact
in \(f^{-1}(W_a)\). Apply \Cref{thm:joint-monotone} on both
spaces to pass each continuous-stage identity to its limit, and sum.
A constant weight gives
\eqref{eq:etale-curvature-degree}.
\end{proof}

\begin{proposition}[Descent of the bounded class]
\label[proposition]{prop:etale-descent}
Let \(f:Y\to X\) be a finite étale surjection and \(u\)
a given real function on \(X\). If
\(u\circ f\in\Cb(Y)\), then \(u\in\Cb(X)=\Rb(X)\).
Membership of a given function in this class is
also local for étale covers.
\end{proposition}
\begin{proof}
Classical subharmonicity descends under the finite surjection
by \cite[Proposition~4.2(6)]{GR}.
A finite surjection is proper and closed, so upper
semicontinuity and local boundedness descend as follows. Since
\(\{u\geq a\}=f(\{u\circ f\geq a\})\), upper
semicontinuity descends.
Cover the finite fiber above \(x\) by finitely many open
neighborhoods on which \(u\circ f\) is bounded, let \(V\)
be their union, and choose a common bound \(M\).
The ordinary open \(X\setminus f(Y\setminus V)\)
contains \(x\), and its entire inverse image lies in \(V\).
Surjectivity gives \(|u|\leq M\) there.

For an étale cover, choose a point above \(x\) and a local
finite-étale representative on compact analytic neighborhoods
\cite[\S11.6.3 and Lemma~11.6.5(a)]{CLD25}. Restrict the target to an ordinary
open in its relative interior. The locally separated
interior formula used in \Cref{lem:local-algebraization}
makes its inverse image ordinary open in the source.
The positive-degree component supplies a finite-étale
surjection near \(x\). The first assertion and ordinary
locality apply. Finally \Cref{thm:local-regularization}
gives the local regularization on \(X\).
\end{proof}

\begin{remark}[Descent data]\label[remark]{rem:tuple-pullback}
\Cref{prop:etale-descent} concerns a function already
given on the base. Functions independently chosen on different
sheets need not descend: constants \(0\) and \(1\) on two
disjoint copies of a ball provide an example.
\end{remark}


\appendix
\section{Curve Laplacians and Borel slicing}
\label{sec:curve-kernel}

The slicing formula in this appendix expresses the first-order
current of \Cref{prop:first-order} as an integral of positive curve
Laplacians. Its measure-theoretic form is needed when the test
coefficient is only Borel.
The measure identities in this section are identities on analytic spaces.
In particular, a Borel coefficient is not replaced by a function on the
image of a moment map. Put $D=\dd$. On a smooth analytic curve we use the
sign for which the Laplacian of a convex function on an oriented annulus
has nonnegative second derivative. With weights $-\log\|e\|$, this is the
sign for which a semipositive model metric has positive curvature of mass
equal to the degree of its line bundle.
Original unit-moment forms are used throughout. The integration theory
of \cite[\S15.1 and \S15.3]{CLD25} restricts on these presentations to
the original one: both integrate their Lagerberg representatives with
the canonical finite-degree calibration. Thus the use below of the
more general skeleton formulation in \cite{CLD25} changes neither the
test forms nor their integrals; compare the original calibration in
\cite[\S3.5]{CLD12}.

\begin{lemma}[Countability of the auxiliary fields]\label[lemma]{lem:gauss-countable}
Let $K/\mathbb Q_p$ be finite. Let $\mathcal G_K$ consist of the fields
obtained from $K$ by finitely many operations
\[
 E\longmapsto \widehat{E(Z_1,\ldots,Z_s)}_{\mathrm{Gauss}},
 \qquad |Z_i|=1.
\]
Every $F\in\mathcal G_K$ has a countable dense subfield. Its valuation
ring is a complete excellent discrete valuation ring of mixed
characteristic, with the same value group as $K$.
Every compact analytic domain admitting a finite affinoid cover in a
good Hausdorff $F$-analytic space is metrizable. Locally in such a space,
ordinary analytic opens admit countable compact
exhaustions and countable families of original smooth scalar functions
that are uniformly dense on each fixed compact support.
\end{lemma}
\begin{proof}
A $\mathbb Q$-span of a $\mathbb Q_p$-basis of $K$ is countable and dense;
the subfield it generates is still countable. If $E_0\subset E$ is a
countable dense subfield, then $E_0(Z_1,\ldots,Z_s)$ is dense for the
unit Gauss valuation in $E(Z_1,\ldots,Z_s)$ and hence in its completion.
The valuation ring after one step is the completion of
$E^\circ[Z_1,\ldots,Z_s]_{(\pi)}$. This is a discrete valuation ring
with uniformizer $\pi$ and residue field
$\widetilde E(Z_1,\ldots,Z_s)$. Its completion is excellent.

An affinoid algebra over a field with a countable dense subfield has a
countable dense subset. Evaluation of the absolute values of these
elements embeds its Berkovich spectrum in a countable product of compact
intervals: uniform approximation determines the value of every other
element. The spectrum is therefore compact metrizable. A finite
affinoid cover gives the assertion for compact domains. An ordinary
open in such a compact neighborhood is locally compact and second
countable, and admits a countable compact exhaustion. Original smooth
partitions of unity and uniform smooth density
\cite[Theorem~15.2.6 and Corollary~15.2.7]{CLD25} give the last assertion.
For a fixed compact set one first chooses a larger compact neighborhood
and a smooth cutoff equal to one on the smaller set, and then takes a
countable dense subset of the resulting separable space of continuous
functions.
\end{proof}

\begin{lemma}[Compatibility of curve Laplacians]\label[lemma]{lem:curve-native}
Let $E$ be a complete nontrivially valued field of characteristic zero,
and let $C$ be an ordinary open in the analytification of a smooth
algebraic $E$-curve. A locally bounded subharmonic function $b$ has a
positive Radon Laplacian $\Delta_C b$, and
\begin{equation}\label{eq:curve-original-test}
 \int_C a\,d(\Delta_C b)=\int_C b\,Da
 \qquad(a\in A_c^{0,0}(C)).
\end{equation}
Both sides use the canonical calibration over $E$. The assertion does not
require the residue field of $E$ to be algebraic over a finite field.
\end{lemma}
\begin{proof}
The open $C$ is strictly analytic. Thuillier's
Theorem~3.4.12 gives its positive Radon Laplacian, and
Theorem~3.4.2 gives, on a relatively compact larger open, a decreasing
net $b_\nu$ of lisse subharmonic functions with limit $b$
\cite{Thuillier}. Here \emph{lisse} means piecewise affine on a finite
semistable skeleton after a finite extension; it does not mean original
CLD smooth. The relation between this subharmonicity and the curve
convention used here is \cite[\S3.3.1--\S3.3.2]{GR}.

For a lisse function, \eqref{eq:curve-original-test} follows from
integration by parts on a common finite skeleton for the function and
the original moment presentation of $a$. Along an annular edge the
formula is the one-variable integration-by-parts formula; the endpoint
terms are the weighted outgoing slopes. These are precisely the atoms
of the lisse Laplacian. The edge weights retain the residue
degrees. A finite extension used to obtain the skeleton can be removed
by the calibrated base-extension integration formula
\cite[\S15.3, Proposition~14.3.6]{CLD25}; both sides have the same
normalization. This is also the lisse comparison underlying
\cite[Proposition~4.2]{Wanner}.

On a compact neighborhood of $\operatorname{supp}a$, a tail of
$b_\nu$ is bounded above, while $b$ bounds it below. Integration of a
decreasing net of continuous functions against a finite Radon measure
commutes with its infimum. To see this, a continuous function strictly
above the infimum eventually dominates the net on the compact set:
the sets on which this domination holds form an increasing open cover.
Compactness gives one index. Taking the infimum of the integrals of
continuous majorants proves the assertion by Radon regularity. Apply
this to the positive and negative parts of the fixed measure $Da$.
The Laplacians converge as distributions by Thuillier's construction;
their masses on a smaller compact set are bounded by testing against
one nonnegative lisse cutoff. Uniform density then permits every
continuous compact test, including $a$. Passing to the limit gives
\eqref{eq:curve-original-test}.
\end{proof}

\begin{lemma}[Borel coarea for a smooth relative form]
\label[lemma]{lem:borel-coarea}
Let $F\in\mathcal G_K$, let $n\geq2$, let $W$ be an ordinary algebraic analytic open
of a smooth $n$-fold, and let
$p:W\to(\mathbb G_m^{n-1})^{\an}$ be induced by a dominant algebraic
morphism. In the specified torus coordinates $z_1,\ldots,z_{n-1}$ put
$t_i=-\log|z_i|$, and let
\[
 \Omega=\omega(t)\bigwedge_{i=1}^{n-1}\dpr t_i\wedge\dpp t_i,
 \qquad 0\leq\omega\in C_c^\infty(\mathbb R^{n-1}).
\]
Thus $\Omega$ is presented by these coordinates and its measure is
supported on their standard Gauss skeleton. Let
$\beta\in A_c^{1,1}(W)$, and let $h$ be a
bounded Borel function on a neighborhood of its support.
Writing $\nu_\Omega=\Omega^{\mathrm{Rad}}$ and $C_y=p^{-1}(y)$, one has
\begin{equation}\label{eq:actual-smooth-coarea}
 \int_W h\,\beta\wedge p^*\Omega
 =\int_{\operatorname{Sk}(\mathbb G_m^{n-1})}
       \left(\int_{C_y}h|_{C_y}\,\beta|_{C_y}\right)d\nu_\Omega(y).
\end{equation}
The inner integral is measurable after completion of $\nu_\Omega$, and
its absolute value is integrable. The fibers over this skeleton are
smooth curves over $\mathscr H(y)$.
\end{lemma}
\begin{proof}
Every torus skeleton point has generic algebraic support. Generic
smoothness in characteristic zero therefore makes $C_y$ a smooth
$\mathscr H(y)$-curve. Since the support of $\beta$ is compact, the
calculation is confined to finitely many compact moment presentations.
The flat locus of $p$ is dense; the hypotheses used in the smooth
coarea theorem \cite[Theorem~16.5.4]{CLD25} are thus satisfied.

For the Borel coefficient in \eqref{eq:actual-smooth-coarea}, refine
the moment presentations by the
base coordinates and choose the convenient skeleton cells of
\cite[Lemma~14.1.5]{CLD25}. On each cell $D\subset W$,
the refined tropicalization $t_D:D\to\overline D$ is a homeomorphism
onto a polyhedral cell. The Borel coefficient used on that cell is
$h\circ t_D^{-1}$. Distinct cells having the same tropical image are
kept separate, with their respective values of $h$. Using relative
interiors, or a disjoint Borel subdivision,
avoids double counting their boundaries.

The integration formula is now a finite sum of ordinary Lebesgue
integrals on these calibrated cells. The base and fiber skeletons are
compatible by \cite[Lemma~14.2.10]{CLD25}. If a full-dimensional source
cell $D$ projects onto a full-dimensional base cell $B$, let $d_D,d_B$
be their calibrated finite degrees. Its fiber calibration is $d_D/d_B$;
this is the degree calculation in the proof of
\cite[Theorem~16.5.4]{CLD25}. Consequently ordinary Borel Fubini on
$\overline D\to\overline B$ gives exactly the corresponding summand
of \eqref{eq:actual-smooth-coarea}, with coefficient $d_D$ on the
source and $d_B(d_D/d_B)$ after integration over the base. Cells of
insufficient base or fiber dimension contribute zero, outside a
finite union of lower-dimensional base cells. That union is
$\nu_\Omega$-null. The same computation applied to absolute values
gives absolute integrability, since $h$ is bounded and the smooth
form has compact support. This proves the formula and its
measurability assertion. The construction evaluates $h$ on $D$, not
on the tropical image without a chosen inverse, and is compatible with further
subdivision and common moment refinements.
When using this formula on a prescribed compact source support, a
coordinate cutoff $\omega$ may be chosen equal to one on a neighborhood
of its compact image under $(-\log|z_1|,\ldots,-\log|z_{n-1}|)\circ p$.
The resulting form $\Omega$ on the base torus has compact support in
a closed polyannulus, so the compactness requirement of the coarea theorem is
retained.
\end{proof}

\begin{theorem}[The curve Laplacian kernel]\label[theorem]{thm:actual-kernel}
In the setting of \Cref{lem:borel-coarea}, let $v$ be locally bounded
and classically plurisubharmonic. After restricting to a relatively
compact neighborhood in $W$, there is a measurable family of positive
Radon measures
\[
 \mu_y=\Delta_{C_y}(v|_{C_y})
 \quad\text{for $\nu_\Omega$-almost every }y
\]
whose integral is a positive Radon measure $\mu_{v,p,\Omega}$ on the
analytic domain $W$. It satisfies
\begin{equation}\label{eq:actual-kernel-smooth}
 \langle D[v],a\,p^*\Omega\rangle
 =\int_W a\,d\mu_{v,p,\Omega}
 =\int\left(\int_{C_y}a|_{C_y}\,d\mu_y\right)d\nu_\Omega(y)
 \quad(a\in A_c^{0,0}(W)).
\end{equation}
For every bounded Borel $h$ on the relevant compact support, the same
measure identity gives
\begin{equation}\label{eq:actual-kernel-borel}
 \int_W ha\,d\mu_{v,p,\Omega}
 =\int\left(\int_{C_y}(ha)|_{C_y}\,d\mu_y\right)d\nu_\Omega(y).
\end{equation}
Scalar multiplication and relations between original test forms are
preserved on common refinements.
\end{theorem}
\begin{proof}
The all-extension quantifier in classical plurisubharmonicity implies
that $v|_{C_y}$ is subharmonic over $\mathscr H(y)$. These fields may
be nondiscrete. \Cref{lem:curve-native} gives their positive
Laplacians. Since $\Omega$ has top degree on the base, it is closed,
and distributional differentiation gives
\[
 L(a):=\langle D[v],a\,p^*\Omega\rangle
      =\int_W v\,Da\wedge p^*\Omega.
\]
Apply \Cref{lem:borel-coarea} with $h=v$ and $\beta=Da$, followed by
\eqref{eq:curve-original-test} on each fiber. This proves that
$y\mapsto\int a|_{C_y}\,d\mu_y$ is measurable modulo a null set and
integrable, and that its integral is $L(a)$. In particular $L(a)\geq0$
when $a\geq0$.

Positivity of $L$ gives an order-zero estimate. Fix
a compact $Q\Subset W$ and an original smooth cutoff $\kappa\geq0$
equal to one near $Q$. For every original smooth $a$ supported in $Q$,
\begin{equation}\label{eq:kernel-mass-bound}
 |L(a)|\leq\|a\|_\infty L(\kappa)<\infty.
\end{equation}
The last value is finite because $v$ is bounded on the compact
support of $D\kappa$ and $D\kappa\wedge p^*\Omega$ is a smooth
top-degree form of finite total variation. Positivity and
\eqref{eq:kernel-mass-bound}, together with original smooth density,
therefore give a positive Radon measure $\mu_{v,p,\Omega}$ by the
Riesz representation theorem.

The preceding argument constructs a measure from its smooth
pairings. To integrate arbitrary Borel functions fiberwise, we also
need a measurable kernel defined outside one common null set.
Put, on $Q'=\operatorname{supp}\kappa$,
$\nu_y=\kappa|_{C_y}\mu_y$, viewed as a finite measure on the compact
metric space $Q'$ and extended by zero off the fiber. Choose a
countable family of original smooth functions uniformly dense in
$C(Q')$. Their products with $\kappa$ have measurable integrals by
the preceding coarea calculation. Together with the test $\kappa$,
this involves only countably many exceptional null sets; remove one
Borel null set containing their union and set $\nu_y=0$ there.
For the remaining parameters the mass $\nu_y(Q')$ is finite.
Uniform approximation consequently proves measurability of
$y\mapsto\int f\,d\nu_y$ for every $f\in C(Q')$.

For an open subset $A$ of $Q'$, continuous functions
$0\leq f_j\uparrow\mathbf1_A$ show that $y\mapsto\nu_y(A)$ is
measurable. The monotone-class theorem then gives this assertion for
every Borel subset of $Q'$, and gives measurable integration of every
bounded Borel function on $Q'$. Moreover,
\[
 \int \nu_y(Q')\,d\nu_\Omega(y)=L(\kappa)<\infty.
\]
Thus $\nu(A)=\int\nu_y(A)\,d\nu_\Omega(y)$ is a finite Borel measure
on the compact metric space $Q'$, hence a Radon measure. On original
smooth tests it agrees with $\kappa\mu_{v,p,\Omega}$; uniform
density identifies these measures. Taking a countable compact
exhaustion and a countable collection of such cutoffs produces a
single conull parameter set and compatible kernels. Equality of the
measures on $W$ proves \eqref{eq:actual-kernel-borel}.

Finally, if two presentations define the same original complementary
form $\eta$, then for every original smooth scalar $f$ their
coefficient measures have integral $\langle D[v],f\eta\rangle$.
They are equal by density and Radon uniqueness. The same argument
applies to every finite relation among test forms and to multiplication
by an original smooth scalar. These relations hold as measure
identities before multiplication by $h$ and remain true afterwards.
This proves compatibility with common refinements and with
Borel coefficients.
\end{proof}

\begin{remark}\label[remark]{rem:kernel-field-scope}
Only the complete-field curve theory is used over $\mathscr H(y)$
in \Cref{thm:actual-kernel}; the envelope theorem is used over
$F\in\mathcal G_K$. The countability needed to select one conull set
is that of a dense family of tests on $W$, not of the residue fields.
The coordinate-presentation
hypothesis on $\Omega$ is essential: an arbitrary original smooth form
on a torus may have a refined skeleton. In the applications below,
$p$ is formed from the units in a decomposable complementary direction;
a coordinate cutoff $\omega=1$ near the projected test support gives
exactly the displayed class of base forms. When $n=1$, the base is a
point and the statement is \Cref{lem:curve-native}. If $p$ is not dominant,
its top-degree base form pulls back to zero by tropical dimension,
and the corresponding first-order pairing is zero.
\end{remark}

\returnlink{sec:first-order}{first-order positivity}

\section{Bounded curve energy and a localized Green estimate}
\label{sec:curve-energy}

The local estimate below transfers differentiation from a continuous
potential to a compactly supported bounded weight. Its total variation
bound then controls the error when the continuous potential varies
uniformly, even though the weight may be discontinuous. It is used in
\Cref{prop:first-order-energy}; the proper-curve identity supplies
the initial case of \Cref{thm:projective-regularization}.

We retain the curve convention of \Cref{lem:curve-native}.
The field in this appendix is complete and nontrivially valued of
characteristic zero. The curves used below are ordinary opens in
smooth algebraic curves; in particular they are strictly analytic.
The operator \(D\) on a lisse function is its weighted graph
Laplacian. Original smooth functions on a common skeleton are
piecewise smooth graph functions, and the same integration-by-parts
formula applies to them.

\begin{lemma}[Bounded curve energy]\label[lemma]{lem:bounded-curve-energy}
Locally bounded subharmonic functions have locally finite Dirichlet
energy. For a finite real linear combination \(b\) of such functions,
the expression
\[
 \Gamma_C(b)=\tfrac12D[b^2]-bD[b]
\]
is a nonnegative Radon measure. Its polarization is bilinear and
satisfies Cauchy--Schwarz on compact supports. The product and chain
rules are
\begin{align}
 D[ab]&=aD[b]+bD[a]+2\Gamma_C(a,b),\label{eq:curve-product}\\
 D[\Phi(b)]&=\Phi'(b)D[b]+\Phi''(b)\Gamma_C(b)
 \label{eq:curve-chain}
\end{align}
whenever the functions are locally bounded and \(\Phi\) is \(C^2\)
on a neighborhood of their relevant range.

On a smooth proper connected curve, if \(r,p\) are bounded
subharmonic metrics on the same line, then their gap \(g=r-p\)
satisfies
\begin{equation}\label{eq:proper-curve-energy}
 E_C(g)=-\int gD[g]\geq0.
\end{equation}
The associated Green pairing is symmetric. If \(E_C(g)=0\), the
function \(g\) is constant pointwise.
\end{lemma}
\begin{proof}
We recall how boundedness puts these statements within the
complete-field curve theory. On a compact domain inside a larger
curve open, Thuillier's harmonic interpolants give lisse
subharmonic functions \(b_S\downarrow b\), indexed by finite sets
of nonrigid points containing the boundary
\cite[Theorem~3.4.2]{Thuillier}. Their values on each such set are
the values of the given function \(b\). At an additional fixed inner boundary
these values can be kept fixed by including its finitely many
points in every \(S\).

Take original smooth cutoffs \(0\leq\rho,\kappa\leq1\), with
\(\kappa=1\) near \(\supp\rho\), on the larger domain. A tail of the
interpolants is bounded by \(M+\epsilon\) on these compact supports
if \(|b|\leq M\); this follows from the increasing open cover
\(\{b_S<M+\epsilon\}\). On a common finite skeleton, the graph
product formula gives, with \(x_S^2=\int\rho^2\Gamma_C(b_S)\),
\begin{equation}\label{eq:curve-local-energy-bound}
 x_S^2\leq (M+\epsilon)^2\|D\kappa\|_{\rm TV}
       +2(M+\epsilon)x_S
                       \left(\int\Gamma_C(\rho)\right)^{1/2}.
\end{equation}
Indeed \(\int_{\supp\rho}Db_S\leq
(M+\epsilon)\|D\kappa\|_{\rm TV}\), by testing against \(\kappa\),
and integration by parts in
\(\int \rho^2\,d'b_S\,d''b_S\) gives the displayed mixed term.
Thus the local Dirichlet energies are bounded.

Subtracting the fixed harmonic extension of the inner boundary
values places the interpolants in the Dirichlet space with zero
boundary values. The finite-skeleton harmonic interpolants are the
orthogonal projections to the corresponding finite-dimensional
Dirichlet subspaces. Their energies increase under refinement;
the bound above puts their limit in the Hilbert completion.
For nested point sets $S\subset S'$, orthogonality gives
\[
 \|b_{S'}-b_S\|_D^2
       =\|b_{S'}\|_D^2-\|b_S\|_D^2
\]
after the fixed harmonic extension has been subtracted. The
bounded increasing energies therefore make this net Cauchy. Its
distributional limit is $b$, so its Dirichlet limit represents the
same potential. This is the construction and
projection argument in \cite[\S3.5, especially Lemma~3.5.3 and the
discussion following it]{Thuillier}. It applies simultaneously to
a finite list of bounded subharmonic functions.

Products of bounded Dirichlet functions and \(C^2\) compositions
satisfy the graph product and chain rules. These rules pass from
the interpolants to the Dirichlet limit: boundedness controls the
multipliers, strong Dirichlet convergence controls the derivative
terms, and the interpolants converge to the given functions at
every point. Green-kernel symmetry
\cite[Remark~3.4.15]{Thuillier} and the energy-integrability
identity \cite[Proposition~3.5.7 and its local variant,
Proposition~3.5.10]{Thuillier} identify the resulting pairings with
integration of these bounded potentials against their
Laplacians. Thus the measure in the statement is
\(\tfrac12D[b^2]-bD[b]\), with the Borel product of the original
representative, and not an independently chosen representative.
The finite-skeleton rules give
\eqref{eq:curve-product}--\eqref{eq:curve-chain}; positivity and
Cauchy--Schwarz pass through the same Dirichlet limit.

For the proper assertion, choose a lisse semipositive reference
metric on the line after replacing it by a positive tensor power
if necessary. More generally a fixed difference of such reference
metrics suffices. Its curvature is a finite combination of atoms
at nonrigid points. The gap of a bounded subharmonic metric from
this reference is bounded. Its Green potential therefore has
finite energy by \cite[Proposition~3.5.7]{Thuillier}.
Polarization gives the symmetric Green pairing and
\eqref{eq:proper-curve-energy}, with the mass of each Chern
measure equal to the degree of its line.

Finally, if \(E_C(g)=0\), Cauchy--Schwarz with the elementary Green
function \(g_{x,y}\) gives \(g(x)=g(y)\) for all nonrigid \(x,y\).
In each local frame the subharmonic weights of \(r\) and of
\(p+g(x)\) therefore agree at all such points. They agree everywhere
by \cite[Lemma~3.4.1]{Thuillier}. Thus the gap is constant at every point.
\end{proof}

\begin{lemma}[Localized variation and Green formula]
\label[lemma]{lem:curve-variation}
Let \(b\) be locally bounded subharmonic on a smooth curve open,
and let \(a\in A_c^{0,0}\). Choose original smooth compact cutoffs
\(0\leq\rho,\kappa\leq1\) with \(\rho=1\) near \(\supp a\) and
\(\kappa=1\) near \(\supp\rho\). Suppose \(|b|\leq M\) on a
neighborhood of these supports. Set
\[
 E(f)=\int d'f\wedge d''f,\qquad A_\kappa=\|D\kappa\|_{\rm TV}.
\]
Then \(D[ab]\) is a signed Radon measure and
\begin{equation}\label{eq:localized-variation}
 \begin{split}
 \|D[ab]\|_{\rm TV}\leq {}&
 M\|a\|_\infty A_\kappa+M\|Da\|_{\rm TV}\\
 &+2M\sqrt{E(a)}
                 \bigl(\sqrt{A_\kappa}+2\sqrt{E(\rho)}\bigr).
 \end{split}
\end{equation}
For continuous subharmonic \(w\),
\begin{equation}\label{eq:localized-green}
 \int ab\,D[w]=\int w\,D[ab].
\end{equation}
If the cutoffs are restrictions of fixed original source functions
in \Cref{thm:actual-kernel}, the bound in
\eqref{eq:localized-variation} is integrable over the slicing base.
\end{lemma}
\begin{proof}
Use a decreasing lisse subharmonic regularization \(b_\nu\) on a
larger open. On the relevant compact set, take a tail with
\(|b_\nu|\leq M+\epsilon\). The estimate
\eqref{eq:curve-local-energy-bound} gives
\[
 x_\nu\leq(M+\epsilon)
              \bigl(\sqrt{A_\kappa}+2\sqrt{E(\rho)}\bigr),
 \qquad x_\nu^2=\int\rho^2\,d'b_\nu\wedge d''b_\nu.
\]
The local mass of \(Db_\nu\) on \(\supp a\) is at most
\((M+\epsilon)A_\kappa\). Apply the graph product formula
\[
 D(ab_\nu)=aDb_\nu+b_\nu Da+
              d'a\wedge d''b_\nu+d'b_\nu\wedge d''a.
\]
The first two terms have the first two bounds of
\eqref{eq:localized-variation}, with \(M+\epsilon\); Cauchy--Schwarz
bounds the last two by \(2\sqrt{E(a)}x_\nu\).
The distributional limit is \(D[ab]\), so the uniform variation
bound makes it a measure with the same estimate. Let
\(\epsilon\downarrow0\).

At fixed \(\nu\), approximate \(w\) by lisse subharmonic functions
on a larger domain. Since \(w\) is continuous, the decreasing net
converges uniformly on the compact support by the finite-cover
form of Dini's argument. The finite-skeleton Green formula
therefore passes to \(w\). Now let \(\nu\) decrease to its limit.
The left side integrates \(ab_\nu\) against the fixed positive
measure \(D[w]\). Decreasing continuous-net integration, applied
after decomposing the sign of \(a\), gives its limit \(ab\).
On the right, distributional convergence and the common total
variation bound imply convergence on continuous compact tests:
approximate such a test uniformly by original smooth functions
and use the variation bound for the error. A cutoff equal to one
on $\supp a$ permits applying this observation to $w$.
This proves
\eqref{eq:localized-green}.

For the base integral, finite moment presentations and compact
coefficient bounds give positive original smooth \((1,1)\)-forms
majorizing both signs of \(D\kappa\) and \(Da\). Their restrictions
to the fibers majorize \(A_\kappa(y)\) and \(\|Da_y\|_{\rm TV}\).
Their integrals over the base are finite by
\Cref{lem:borel-coarea}. The same is true of \(E(a_y)\) and
\(E(\rho_y)\), since these are fiber integrals of fixed positive
smooth forms. Cauchy--Schwarz on the base bounds the products of
square roots in \eqref{eq:localized-variation}. This proves
integrability without a uniform degree assertion for the fibers.
\end{proof}

\begin{corollary}\label[corollary]{cor:weighted-uniform-curve}
If continuous subharmonic \(w_j\to w\) uniformly on the relevant
support, then
\[
 \left|\int ab\,(D[w_j]-D[w])\right|
 \leq\|w_j-w\|_\infty\|D[ab]\|_{\rm TV}.
\]
The same conclusion holds for finite real linear combinations of
the bounded subharmonic weight \(b\).
\end{corollary}
\begin{proof}
Subtract \eqref{eq:localized-green} for \(w_j\) and \(w\), and
use the variation norm. Then use linearity.
\end{proof}

\begin{proposition}[The curve contact theorem]\label[proposition]{prop:curve-contact}
Let \(C/F\) be smooth projective geometrically integral, and \(L\)
ample. Let \(p\) be a continuous model-semipositive metric on \(L\)
and \(r\) a bounded classical metric on \(L\). If
\[
 g=r-p\geq0,\qquad \int g\,c_1(L,p)=0,
\]
then \(g=0\) at every point.
\end{proposition}
\begin{proof}
Both local weights are subharmonic. By
\Cref{lem:bounded-curve-energy},
\[
 0\leq E_C(g)=-\int gD[g]
       =\int g\,c_1(L,p)-\int g\,c_1(L,r)\leq0.
\]
Hence \(g\) is constant. The measure \(c_1(L,p)\) has
positive mass \(\deg_F L\), so the constant is zero.
\end{proof}
\returnlink{prop:first-order-energy}{smooth-background energy}
\returnlink{thm:projective-regularization}{projective contact and regularization}

\section{Envelopes on a normal model}
\label{sec:normal-envelopes}

We prove continuity of the ample envelopes used in
\Cref{thm:projective-regularization}. Weighted base ideals encode the
semipositive minorants of a model obstacle. We replace them by test
ideals to obtain global generation after one fixed ample twist.
The required comparison comes from subadditivity with a vertical
correction independent of the ideals and their exponents. Dividing
logarithmic norms by degree makes the resulting error tend to zero.

Throughout this appendix $R$ is a complete discrete valuation ring of
mixed characteristic, $F=\operatorname{Frac}(R)$ has characteristic
zero, $\pi$ is a uniformizer, and $\lambda=-\log|\pi|>0$. No finiteness
or perfection assumption is imposed on the residue field. All test
ideals are those of \cite[Definition~8.37]{BMP}; after completion we use
the stable complete-local ideals of \cite[Definition~8.13]{BMP}.

\subsection{A correction independent of ideals and exponents}

Let $(S,\mathfrak m)$ be a complete Noetherian normal local domain
of mixed characteristic $(0,p)$ and dimension $d$, and let
$\Delta\geq0$ be an integral Weil divisor such that
$K_S+\Delta$ is Cartier. Choose a local trivialization of this Cartier
divisor and a compatible rational canonical section. The canonical
module is then represented by
\[
 W=\mathcal O_S(-\Delta)\subset S,
 \qquad K_S=-\Delta.
\]
Choose $c\ne0$ with $cS\subset W$. Put
\[
 E=H^d_{\mathfrak m}(W),\quad F_S=H^d_{\mathfrak m}(S),\quad
 \alpha:E\longrightarrow F_S,\quad \beta:F_S\longrightarrow E,
\]
where $\alpha$ is induced by inclusion and $\beta$ by multiplication
by $c$. Both composites are multiplication by $c$. Top local duality
identifies the Matlis dual of $E$ with
$\operatorname{Hom}_S(W,W)=S$, so $E$ is the injective hull of the
residue field. This does not require $S$ to be Cohen--Macaulay, and
$\alpha$ need not be injective.

The Noether different will supply the local correction factor. Its
role is to factor multiplication on a big Cohen--Macaulay algebra
through a flat module, where annihilators of finitely generated
ideals commute with base change.

\begin{lemma}[A flat factorization]\label[lemma]{lem:different-factor}
Let $A\subset S$ be a finite extension of complete local domains, with
$A$ regular of dimension $d$. Write
\[
 \mathfrak N_{S/A}
 =\mu\bigl(\operatorname{Ann}_{S\otimes_A S}(\ker\mu)\bigr),
 \qquad \mu:S\otimes_A S\longrightarrow S.
\]
If $B$ is a balanced big Cohen--Macaulay $S^+$-algebra, multiplication
by every $z\in\mathfrak N_{S/A}$ on $B$ factors $S$-linearly as
\[
 B\xrightarrow{j_z}S\otimes_A B\xrightarrow{\mu_B}B,
 \qquad \mu_Bj_z=z\operatorname{id}_B,
\]
and $S\otimes_A B$ is $S$-flat. For every finitely generated ideal
$I\subset S$, it follows that
\begin{equation}\label{eq:different-annihilator}
 z\operatorname{Ann}_{H^d_{\mathfrak m}(B)}(I)
 \subset
 \operatorname{im}\bigl((S\otimes_A B)\otimes_S
       \operatorname{Ann}_{F_S}(I)\longrightarrow H^d_{\mathfrak m}(B)\bigr).
\end{equation}
\end{lemma}
\begin{proof}
Every system of parameters of $A$ is a system of parameters of $S$.
Thus $B$ is balanced big Cohen--Macaulay over $A$, and is $A$-flat by
\cite[Setup~2.1 and Proposition~4.9]{Holm}. Choose
$\gamma=\sum u_i\otimes v_i$ annihilating $\ker\mu$ and satisfying
$\mu(\gamma)=z$. Set $j_z(b)=\sum u_i\otimes v_i b$.
The relations $(s\otimes1-1\otimes s)\gamma=0$ prove $S$-linearity;
multiplication proves the asserted composite. Flat base change gives
the flatness of the middle module $C=S\otimes_A B$.
A top \v{C}ech cokernel gives $H^d_{\mathfrak m}(C)=C\otimes_S F_S$.
Since $I$ is finitely generated and $C$ is flat,
$\operatorname{Ann}_{C\otimes F_S}(I)=C\otimes\operatorname{Ann}_{F_S}(I)$.
Apply local cohomology to the factorization and then apply $\mu_B$.
\end{proof}

\begin{proposition}[Complete-local corrected subadditivity]
\label[proposition]{prop:local-subadditivity}
With the preceding notation, for nonzero ideals $\mathfrak a,\mathfrak b$
and rational $s,t\geq0$,
\begin{equation}\label{eq:local-corrected}
 c\mathfrak N_{S/A}\,
 \tau(S,\Delta,\mathfrak a^s\mathfrak b^t)
 \subset
 \tau(S,\Delta,\mathfrak a^s)\tau(S,\Delta,\mathfrak b^t).
\end{equation}
\end{proposition}
\begin{proof}
Choose a single sufficiently large perfectoid balanced BCM
$S^+$-algebra $B$ for the three stable test ideals. Such a choice, and
the mixed element-wise description below, are supplied by
\cite[Definitions~8.2 and~8.13, Theorem~8.11]{BMP}.
For an ideal datum $\mathfrak a^s$, let $\mathcal R_{\mathfrak a,s}$
be the allowed root multipliers: powers $f^q$ in $B$, where
$0\ne f\in\mathfrak a^m$ and $q\geq s/m$ is rational. Put
$T_{\mathfrak a}=\tau(S,\Delta,\mathfrak a^s)$. Then
\begin{equation}\label{eq:pair-kernel}
 \operatorname{Ann}_E(T_{\mathfrak a})
 =\bigcap_{g\in\mathcal R_{\mathfrak a,s}}
       \ker\bigl(E\xrightarrow{\alpha}F_S
                      \longrightarrow H^d_{\mathfrak m}(B)
                      \xrightarrow{g}H^d_{\mathfrak m}(B)\bigr).
\end{equation}
Indeed, the root maps defining the parameter test modules start from
$F_S$, whose Matlis dual is $W$. The dual of $\alpha$ is the chosen
inclusion $W\hookrightarrow S=\operatorname{Hom}_S(W,W)$.
Precomposition therefore gives precisely the pair ideals as ideals of
$S$. The equality follows from Matlis duality. The pair factor is
one because $K_S+\Delta=0$ in the chosen trivialization.

Here stability is required simultaneously for the three ideal data,
rather than for every ideal on $S$ at once. The existence and
comparison theorem in \cite[Theorem~8.11]{BMP} permits replacing
their three sufficiently large algebras by a common one. The mixed
version of \cite[Definition~8.2]{BMP} is a sum of the element-wise
parameter test modules for products of root multipliers. Matlis
duality changes that sum into the intersection of kernels in
\eqref{eq:pair-kernel}, and gives the corresponding intersection
over pairs $(g,h)$ for the mixed datum.

Take $\eta\in\operatorname{Ann}_E(T_{\mathfrak a}T_{\mathfrak b})$
and $h\in\mathcal R_{\mathfrak b,t}$. For every $a\in T_{\mathfrak a}$,
the element $a\eta$ is annihilated by $T_{\mathfrak b}$; its image
under $h\alpha$ is therefore zero by \eqref{eq:pair-kernel} for
$\mathfrak b^t$. Hence $h\alpha(\eta)$ is annihilated by
$T_{\mathfrak a}$.
For $z\in\mathfrak N_{S/A}$, \eqref{eq:different-annihilator} expresses
\[
 zh\alpha(\eta)=\sum_i b_i\zeta_i,
 \qquad \zeta_i\in\operatorname{Ann}_{F_S}(T_{\mathfrak a}),\quad b_i\in B.
\]
Now $\beta(\zeta_i)\in\operatorname{Ann}_E(T_{\mathfrak a})$.
Equation~\eqref{eq:pair-kernel} gives
$gc\zeta_i=g\alpha\beta(\zeta_i)=0$ for every
$g\in\mathcal R_{\mathfrak a,s}$. Hence $czgh\alpha(\eta)=0$.
The mixed element-wise definition tests exactly the products $gh$.
Thus
\[
 cz\operatorname{Ann}_E(T_{\mathfrak a}T_{\mathfrak b})
 \subset\operatorname{Ann}_E
       \tau(S,\Delta,\mathfrak a^s\mathfrak b^t).
\]
Taking annihilators yields \eqref{eq:local-corrected}. No factorization
of an arbitrary root of an element of $(\mathfrak a\mathfrak b)^m$
is used.
\end{proof}

To globalize the correction, we need finitely many local differents
whose sum is the unit ideal on the generic fiber. The next lemma
obtains them from finite projections that are \'etale on a covering
of that fiber. Their common failure is then supported vertically and
is killed by one power of $\pi$.

\begin{lemma}[Finite projections and a vertical different]
\label[lemma]{lem:finite-projections}
Let $\mathcal X$ be normal, integral, projective and flat over $R$,
with smooth generic fiber $X$ of dimension $n$. There are finitely
many finite morphisms $f_j:\mathcal X\to\mathbb P_R^n$ whose generic
\'etale loci cover $X$. The sum $\mathfrak J$ of their coherent
Noether-different ideals contains $\pi^{C_0}\mathcal O_{\mathcal X}$
for some $C_0$.
\end{lemma}
\begin{proof}
Choose a relatively very ample line $\mathcal H$. On the special
fiber, choose $n+1$ homogeneous parameters without a common zero.
This works over an arbitrary residue field: at each stage, Serre
interpolation provides a high-degree section nonzero at a chosen
closed point of each component of the preceding zero locus. The
dimension drops at each step. Taking powers makes the degrees equal;
taking further powers permits lifting the sections to $\mathcal X$
by Serre vanishing. Every tuple congruent to these lifts modulo $\pi$
has no base point on $\mathcal X$: a nonempty proper closed base locus
would meet the special fiber. It defines a finite morphism to
$\mathbb P_R^n$, because the pullback of $\mathcal O(1)$ is relatively
ample and so no fiber can have positive dimension.

Fix a closed point $x\in X$. Increase the degree until sections
surject onto its first infinitesimal neighborhood. Since
$\kappa(x)/F$ is separable, one can prescribe one nonzero section and
$n$ ratios with independent differentials at $x$. Being \'etale at
$x$ is therefore a nonempty Zariski-open condition on the tuple; the
determinant condition can be expressed over $F$ by its field norm.
The congruence coset of lifts modulo $\pi$ is Zariski dense in the
$F$-vector space of tuples, since $F$ is infinite and the section
module is a full $R$-lattice. Choose a tuple in this open set.
The resulting generic \'etale loci cover every closed point of $X$,
hence all of the Jacobson scheme $X$. Quasi-compactness selects
finitely many morphisms.

On affine target charts their Noether-different ideals glue, since
the diagonal ideal is finitely generated and annihilators commute
with localization. On the smooth generic fiber the different is a
unit at every \'etale point; see \cite[\S49.12]{Stacks}.
Thus $\mathfrak J|_X=\mathcal O_X$. The coherent quotient
$\mathcal O_{\mathcal X}/\mathfrak J$ is supported on the special
fiber. A finite affine cover gives one exponent $C_0$.
\end{proof}

\begin{theorem}[Global fixed correction]\label[theorem]{thm:fixed-correction}
Let $\mathcal X/R$ be as in \Cref{lem:finite-projections}. Fix an
effective integral vertical Weil divisor $\Delta_0$ such that
$K_{\mathcal X}+\Delta_0$ is Cartier. There is an integer $C\geq0$,
independent of the nonzero coherent ideals $\mathfrak a,\mathfrak b$
and rational exponents $s,t\geq0$, such that
\begin{equation}\label{eq:global-correction}
 \pi^C\tau(\mathcal X,\Delta_0,\mathfrak a^s\mathfrak b^t)
 \subset
 \tau(\mathcal X,\Delta_0,\mathfrak a^s)
 \tau(\mathcal X,\Delta_0,\mathfrak b^t).
\end{equation}
\end{theorem}
\begin{proof}
Choose $T$ globally with
$\pi^T\mathcal O_{\mathcal X}\subset\mathcal O_{\mathcal X}(-\Delta_0)$.
It exists because the latter coherent ideal is the unit ideal on
$X$. It gives $c=\pi^T$ in every local canonical embedding above.

Fix a closed point $y$ of the special fiber and a finite projection
from \Cref{lem:finite-projections}. Complete the regular local ring
of $\mathbb P_R^n$ at the image of $y$. Its finite algebra splits
into complete local factors; the factor at $y$ is
$S=\widehat{\mathcal O}_{\mathcal X,y}$. It is a normal domain by
excellence. Both rings have dimension $n+1$, so the map from the
complete regular domain $A$ to $S$ is injective. An original
Noether-different tensor witness base changes to the finite completed
algebra and projects in both tensor factors to a witness for $S/A$.
Thus the completed image of the global different is contained in
$\mathfrak N_{S/A}$.

Apply \Cref{prop:local-subadditivity}. Corollary~8.40 of \cite{BMP}
identifies the completed global test ideals with the complete-local
ones, and faithful flatness descends the containment. Sum over the
finitely many projections and use
$\pi^{C_0}\mathcal O\subset\mathfrak J$. This gives
\eqref{eq:global-correction} with $C=T+C_0$ at every closed
special-fiber point. A nonzero coherent failure quotient on the proper
$R$-scheme would have such a point in its support. The containment
therefore holds everywhere, with a constant chosen before the ideals
and exponents.
\end{proof}

\begin{remark}\label[remark]{rem:regular-subadditivity}
The uncorrected subadditivity in \cite[Theorem~8.41(h)]{BMP} applies
on a regular scheme without a boundary in its two factors.
In \Cref{thm:fixed-correction}, the scheme is normal and the
boundary $\Delta_0$ occurs in both factors. The factor $\pi^C$
accounts for these differences. Its independence of the ideals and
the exponents will give an error of order $m^{-1}$ in the envelope
construction.
Fang--Gubler--K\"unnemann prove continuity for model obstacles on
regular models carrying an ample extension, and for arbitrary
continuous obstacles under resolution of projective models
\cite[Theorem~1.1]{FGK}. The proof below instead uses
\Cref{thm:fixed-correction} on one normal model and retains a
uniform error term throughout the construction.
\end{remark}

\subsection{Weighted ideals and removal of the auxiliary line}

\begin{theorem}[Continuity and orthogonality of ample envelopes]
\label[theorem]{thm:normal-envelope}
Let $X/F$ be smooth projective geometrically integral and $L$ ample.
For every continuous metric $\rho$ on $L$, its pointwise envelope
\[
 P_L\rho=\sup\{q:q\text{ is continuous model-semipositive on }L, q\leq\rho\}
\]
is continuous and model-semipositive on the same identified line.
It is uniformly approximable by nef rational model metrics over $F$,
and
\begin{equation}\label{eq:envelope-orthogonality}
 \int_{X^{\an}}(\rho-P_L\rho)c_1(L,P_L\rho)^n=0.
\end{equation}
\end{theorem}
\begin{proof}
Take a normal projective flat model $\mathcal X$ carrying an extension
$\mathcal L$ of $L$. Replace it by the normalization of the blowup
of its rank-one canonical fractional sheaf. This modification is an
isomorphism on the smooth
generic fiber and gives a Cartier extension $G$ of its canonical line.
Compatible rational sections make $G-K_{\mathcal X}$ vertical.
For sufficiently large $N_0$,
\[
 \Delta_0=G+N_0\operatorname{div}(\pi)-K_{\mathcal X}\geq0,
 \qquad K_{\mathcal X}+\Delta_0\text{ Cartier}.
\]
First suppose that the obstacle is an integral model metric $r$;
rational denominators will be cleared and divided out afterwards.
Let $\theta$ be the metric of $\mathcal L$. On a dominating normal
model $f:\mathcal W\to\mathcal X$, a constant translation permits
writing
\[
 r=\theta+u,\qquad u=\log|\mathcal E|,
 \qquad (\pi^N)\mathcal O_{\mathcal W}\subset\mathcal E
                  \subset\mathcal O_{\mathcal W},
\]
with $\mathcal E$ an invertible vertical ideal. The section lattices
below encode the condition that a Fubini--Study weight lie below $r$.
We form their base ideals $\mathfrak c_l$ and the associated
asymptotic test ideals $\beta_m$:
\begin{align}
 V_l&=H^0\bigl(\mathcal W,(f^*\mathcal L\otimes\mathcal E)^l\bigr),\notag\\
 \mathfrak c_l&=\operatorname{im}
       (V_l\otimes\mathcal L^{-l}\longrightarrow\mathcal O_{\mathcal X}),
       \label{eq:weighted-ideals}\\
 \beta_m&=\sum_{l\gg0}\tau(\mathcal X,\Delta_0,\mathfrak c_l^{m/l}),
 \quad \psi_m=m^{-1}\log|\beta_m|,
 \quad P=P_Lr-\theta.\notag
\end{align}
Only sufficiently large degrees are used, so these ideals are vertical.
The sum is directed under divisibility, since
$\mathfrak c_l^q\subset\mathfrak c_{lq}$ and test ideals commute
with powers. It stabilizes by Noetherianity.

Normality identifies $V_l$ with the unit lattice for the sup norm of
$lr$. Pure-FS density on an ample line
\cite[Corollary~7.6 and Theorem~7.8]{BE} gives
\begin{equation}\label{eq:weighted-envelope-sup}
 P=\sup_{l\gg0}l^{-1}\log|\mathfrak c_l|,
 \qquad (mq)^{-1}\log|\mathfrak c_{mq}|\longrightarrow P
 \quad(q\longrightarrow\infty)
\end{equation}
for fixed $m$. For the latter limit, products of unit sections give
superadditivity of the degree-weighted functions. Divide large
degrees by any fixed sufficiently large $l$, leaving remainders in
a finite interval of globally generated degrees. Their weights are
bounded on the compact analytic space, so their contribution divided
by the total degree tends to zero. Taking the supremum over $l$
proves the stated pointwise limit.

Let $\tau_0=\tau(\mathcal X,\Delta_0)$. It contains a power of $\pi$,
and Cartier translation and monotonicity give
$\mathfrak c_{mq}\tau_0\subset\beta_{mq}$. Iterating
\eqref{eq:global-correction} and using repeated-factor compatibility
\cite[Theorem~8.41(c)]{BMP}, choose a degree computing $\beta_{mq}$ to get
\begin{equation}\label{eq:asymptotic-correction}
 \pi^{C(q-1)}\beta_{mq}\subset\beta_m^q.
\end{equation}
Indeed, \eqref{eq:asymptotic-correction} gives
\[
 \frac{1}{mq}\log|\beta_{mq}|
       -\frac{C(q-1)\lambda}{mq}\leq\psi_m.
\]
Use $\mathfrak c_{mq}\tau_0\subset\beta_{mq}$ and
\eqref{eq:weighted-envelope-sup}, then let $q\to\infty$.
The contribution of the fixed ideal $\tau_0$ tends to zero, leaving
$P-C\lambda/m\leq\psi_m$.

There is also a fixed upper bound
\begin{equation}\label{eq:ideal-upper}
 \psi_m\leq u+B\lambda/m.
\end{equation}
Indeed, put
$\Delta_{\mathcal W}=f^*(K_{\mathcal X}+\Delta_0)-K_{\mathcal W}$.
For a noneffective $\Delta_{\mathcal W}$, choose an effective
vertical Cartier divisor $A$ with $\Delta_{\mathcal W}+A\geq0$ and
interpret its test ideal as
$\mathcal O(A)\tau(\mathcal W,\Delta_{\mathcal W}+A)$.
Cartier translation makes this independent of $A$.
Birational inclusion, monotonicity, powers and Cartier translation
\cite[Proposition~7.16(b), Theorem~8.41(b),(c),(j)]{BMP} imply
\[
 \tau(\mathcal X,\Delta_0,\mathfrak c_l^{m/l})\mathcal O_{\mathcal W}
 \subset\mathcal E^m\tau(\mathcal W,\Delta_{\mathcal W}).
\]
Indeed, the pushforward inclusion of
\cite[Theorem~8.41(j)]{BMP} gives the first pullback containment
by adjunction. The further containment follows from
$\mathfrak c_l\mathcal O_{\mathcal W}\subset\mathcal E^l$ and
Cartier translation for the invertible ideal $\mathcal E$.
The last test ideal is a fixed coherent fractional ideal, equal to
$\mathcal O_X$ on the generic fiber. One power $\pi^B$ clears its
vertical poles, proving \eqref{eq:ideal-upper}.

The bounds on $\psi_m$ compare functions, but do not yet give
semipositive metrics. For that purpose we show that one fixed ample
twist makes every $\beta_m\mathcal L^m$ globally generated. Fix $l$ and choose $R$-module generators
$\sigma_1,\ldots,\sigma_a$ of $V_l$, regarded as sections of
$\mathcal L^l$ on $\mathcal X$. Discard any zero generator and let
$D_b=\operatorname{div}(\sigma_b)$, an effective Cartier divisor
linearly equivalent to $l\mathcal L$. Locally trivializing
$\mathcal L^l$ expresses $\mathfrak c_l$ as the sum of the
principal ideals defined by these sections. Repeated application of
the summation formula \cite[Theorem~8.41(g)]{BMP}, together with
Cartier translation, gives
\[
 \tau(\mathcal X,\Delta_0,\mathfrak c_l^{m/l})
 =\sum_{\substack{t_b\in\mathbb Q_{\geq0}\\\sum_b t_b=m/l}}
   \tau\left(\mathcal X,\Delta_0+\sum_b t_bD_b\right).
\]
The local expressions agree on overlaps because they use the
Cartier divisors of global sections. Each divisor
$D=\sum_b t_bD_b$ is effective and satisfies
$D\sim_{\mathbb Q}m\mathcal L$. By Noetherianity, finitely many
terms suffice for this sum and for the sum defining $\beta_m$;
compare \cite[Remark~8.42]{BMP}.

Choose a generated ample Cartier divisor $H$ and write
$G_0=K_{\mathcal X}+\Delta_0$. This is Cartier by construction.
For each such $D$, apply \cite[Corollary~7.24]{BMP} with
$M=m\mathcal L+G_0+H$. Its required difference is
\[
 M-K_{\mathcal X}-(\Delta_0+D)
       \sim_{\mathbb Q}H,
\]
which is big and nef. Since the special fiber has dimension $n$,
the sheaf
\[
 \tau(\mathcal X,\Delta_0+D)
       \otimes\mathcal L^m\otimes\mathcal O(G_0+(n+1)H)
\]
is generated. Take a generated ample model line $\mathcal A$ and
choose $m_0$, independently of $m,l,D$, so that
$\mathcal A^{m_0}\otimes\mathcal O(-G_0-(n+1)H)$ is generated.
Tensoring with this fixed generated line and then taking the sum
proves that $\beta_m\mathcal L^m\mathcal A^{m_0}$ is generated for
every $m$.

Let $A=\mathcal A|_X$ with its model metric $q_A$. Generation gives
a pure-FS metric $t_m$ on $mL+m_0A$ with exact weight
$m\theta+m\psi_m+m_0q_A$. This metric is on $mL+m_0A$;
to compare it with the envelope, we must return to $L$.
Choose $r_0$ so that $r_0L-A$ is generated, and a pure-FS metric
$\nu$ on that line. Adding $m_0\nu$ cancels the auxiliary $A$
and produces a metric on $(m+r_0m_0)L$. Set
\[
 M_m=m+r_0m_0,\qquad h_m=(t_m+m_0\nu)/M_m,\qquad
 d=q_A+\nu-r_0\theta,\quad D_0=\|d\|_\infty.
\]
Every $h_m$ is a nef rational model metric on $L$. Since
$-N\lambda\leq u\leq0$ and $P\leq u$, the lower estimate and
\eqref{eq:ideal-upper} give, with a constant $C_1$ independent of $m$,
\[
 h_m-P_Lr\geq-\frac{C\lambda+m_0D_0}{M_m},
 \qquad h_m-r\leq\frac{C_1}{M_m}.
\]
Shift $h_m$ down by its upper error $C_1/M_m$. The result is a
continuous semipositive minorant of $r$, hence lies below $P_Lr$
by the definition of the envelope. Combining this with the lower
estimate gives
\begin{equation}\label{eq:envelope-uniform-bound}
 -\frac{C\lambda+m_0D_0}{M_m}
 \leq h_m-P_Lr\leq\frac{C_1}{M_m}.
\end{equation}
This proves uniform convergence for a model obstacle on the prescribed
line. Real constant translations are allowed in the uniform closure
of rational model metrics; rational downward approximations give the
same pointwise envelope. Undo the preparatory translation.

For a rational model obstacle clear its tensor denominator, use
$P_{dL}(dr)=dP_Lr$, and divide back. For a continuous obstacle, uniform
model density and the estimate
$\|P_Lr-P_Lr'\|_\infty\leq\|r-r'\|_\infty$
give continuity and same-line semipositive approximation by a diagonal.
Finally \cite[Theorem~6.7]{BGM}, applicable because $X$ is geometrically
reduced, gives \eqref{eq:envelope-orthogonality} after this continuity
has been established.
\end{proof}

\begin{remark}\label[remark]{rem:gauss-envelope-scope}
Every field in $\mathcal G_K$ satisfies the hypotheses of this
appendix by \Cref{lem:gauss-countable}. The proof applies to arbitrary
projective varieties, lines and obstacles defined over that field;
it does not require their descent to $K$. This is the field scope
needed for the induction through generic hyperplanes.
\end{remark}
\returnlink{thm:projective-regularization}{projective contact and regularization}

\section{Algebraic coordinates and a common local domain}
\label{sec:local-geometry}

The projective approximation theorem must be applied to local
potentials without changing them near the test support. The collar
in \Cref{lem:common-collar} achieves this for an entire decreasing
family at once: its first terms and limits determine a common bound,
and that bound determines one barrier and one local domain.
The last two lemmas transfer these constructions to the analytic
spaces of \Cref{conv:spaces}, including their integration
normalizations.

In the first three lemmas, \(F\in\mathcal G_K\),
\(\pi\) is the chosen uniformizer, and \(\lambda=-\log|\pi|\).
All constructions take place over \(F\) on the indicated ordinary
neighborhoods.

\begin{lemma}[Exact replacement of analytic units]
\label[lemma]{lem:exact-units}
Let \(Y/F\) be smooth integral affine, let \(O\subset Y^{\an}\) be ordinary open,
and let \(f_1,\ldots,f_a\) be analytic units on \(O\).
Every \(x\in O\) has a compact affinoid neighborhood \(C\subset O\)
and rational functions \(g_1,\ldots,g_a\in F(Y)\) such that
\begin{equation}\label{eq:exact-units}
 \|f_i-g_i\|_C<\min_C|f_i|,
 \qquad |g_i(y)|=|f_i(y)|\quad(y\in C).
\end{equation}
The \(g_i\) are units on one algebraic open containing \(C\).
Consequently the logarithms and their original smooth differentials
agree on the ordinary interior of \(C\). A compactly supported
smooth test can be divided into finitely many tests on such
neighborhoods, with one simultaneous replacement for each finite
moment presentation.
\end{lemma}
\begin{proof}
Embed \(Y\) in affine space with
coordinates \(T_1,\ldots,T_N\).
Polynomial evaluations generate the topology of \(Y^{\an}\).
Since \(|F^\times|^{\mathbb Q}\) is dense in \(\mathbb R_{>0}\),
there are finitely many polynomial inequalities, with strict slack
at \(x\), whose closed solution set \(C\) contains an ordinary
neighborhood of \(x\) and lies in \(O\).
Include upper bounds for every \(T_j\). Thus \(C\) is compact, and
it has a presentation
\[
 |T_j|\le R_j,\qquad |P_i|\le S_i,\qquad |Q_l|\ge r_l>0,
 \qquad R_j,S_i,r_l\in |F^\times|^{\mathbb Q}.
\]
It is obtained from the closed polydisc by imposing the equations of
\(Y\), followed by Weierstrass and Laurent localizations.
Its affinoid algebra is a quotient of the generalized Tate algebra
with coordinates \(T_j,U_i,V_l\) and relations
\[
 U_i=P_i(T),\qquad Q_l(T)V_l=1.
\]
The radii for \(U_i\) and \(V_l\) are \(S_i\) and \(r_l^{-1}\).
Polynomials in these coordinates are dense in that Tate algebra,
and hence their images are uniformly dense in \(\cO(C)\).
Those images belong to
\(F[Y][Q_1^{-1},\ldots,Q_b^{-1}]\).
Compactness and invertibility give \(\min_C|f_i|>0\).
Approximate the finitely many \(f_i|_C\) by elements of the displayed
localization with errors smaller than these minima.
The ultrametric inequality gives the pointwise equalities in
\eqref{eq:exact-units}. Their numerators and denominators do not vanish
anywhere on \(C\); removing their algebraic zero loci gives one
algebraic open on which all \(g_i\) are units.
The logarithms are equal as smooth functions on \(\operatorname{Int}C\),
so applying \(\dpr,\dpp\) preserves the equality.
For the last assertion, cover the compact test support by finitely
many such interiors and use a smooth partition subordinate to this
finite cover. Apply the construction to the whole finite unit list
on each summand.
\end{proof}

\begin{lemma}[A common collar for decreasing families]
\label[lemma]{lem:common-collar}
Let \(Y/F\) be smooth projective and let \(O\subset Y^{\an}\) be
ordinary open. For finitely many indices \(i\), let
\(u_{i,j}\downarrow u_i\) be classical plurisubharmonic functions
whose first terms and limits are locally bounded.
Near every \(x\in O\), there are one ordinary neighborhood \(V\)
with compact closure in \(O\), an ample line bundle \(L\) on \(Y\),
an algebraic frame \(e\) near \(\overline V\), a constant
\(c\in\lambda\mathbb Q\), and bounded classical metrics
\(R_{i,j}\downarrow R_i\) on this same \(L\) such that
\begin{equation}\label{eq:common-collar-weights}
 \operatorname{wt}_e(R_{i,j})=u_{i,j}+c,\qquad
 \operatorname{wt}_e(R_i)=u_i+c\quad\hbox{on }V.
\end{equation}
All geometric choices are independent of \(i,j\).
\end{lemma}
\begin{proof}
Choose a projective embedding by a sufficiently ample line \(A\)
and an affine coordinate chart containing \(x\).
As in the proof of \Cref{lem:exact-units}, choose nested polynomial
norm boxes in this chart,
\[
 x\in\operatorname{Int}C_1,\qquad
 C_1\subset\operatorname{Int}C_{\mathrm{mid}},\qquad
 C_{\mathrm{mid}}\subset\operatorname{Int}C_2,\qquad C_2\subset O,
\]
using the same finite polynomial list and strictly nested thresholds.
All lower-bound polynomials are nonzero on \(C_2\).
Choose additional thresholds strictly between those for \(C_1\)
and \(C_{\mathrm{mid}}\).
Clear rational powers in the thresholds to obtain a function
\begin{equation}\label{eq:collar-barrier}
 b=\max\left\{0,\log|P_a^{N_a}/c_a|,
                    \log|d_l/Q_l^{M_l}|\right\},
 \qquad c_a,d_l\in F^\times .
\end{equation}
The barrier must vanish on the inner box, where we want to retain
the given potentials, and be uniformly positive on the outer collar,
where the extension must agree with a global metric. The function
$b$ is finite continuous on $C_2$, zero on $C_1$, and satisfies
\[
 b\ge\delta>0\quad\hbox{on }C_2\setminus\operatorname{Int}C_{\mathrm{mid}}
\]
for one \(\delta\).
Indeed, outside the strict intermediate box at least one defining
bound fails; the ratio between that threshold and its chosen inner
threshold is strictly larger than one. The minimum of the finitely
many positive logarithmic gaps gives \(\delta\).

Clear the denominators in \eqref{eq:collar-barrier} and homogenize.
After increasing the common degree by a power of the affine-chart
denominator, the branches become ratios \(s_a/s_0\) of sections of
\(A^{m_0}\), where \(s_0\) is nowhere zero on \(C_2\); include
the branch \(s_0/s_0=1\). This denominator can be chosen as a product
of the homogenized \(Q_l\) and the affine-chart denominator, so its
nonvanishing on \(C_2\) is preserved.
Choose \(m_0\) large enough that \(A^{m_0}\) is generated.
Append small \(F\)-scalar multiples of a global generating tuple.
Since \(C_2\) is compact and \(s_0\) has no zero there, the appended
ratios have absolute value less than one on \(C_2\).
Their ratios stay below the branch $1$ on $C_2$, so they do not
change $b$ there. They ensure that the enlarged tuple generates the
line everywhere. Its pure Fubini--Study metric \(q\) has
weight \(b\) in the frame \(s_0\) on \(C_2\).

The inequalities \(u_i\le u_{i,j}\le u_{i,1}\) and compactness give
one two-sided bound for all the functions on \(C_2\).
Choose \(c\in\lambda\mathbb Q\) and \(M\) with
\(0<u_{i,j}+c\le M\) and \(0<u_i+c\le M\) there.
Choose an integer \(m>0\) with \(m\delta>M+1\), and put
\(L=A^{mm_0}\), \(e=s_0^m\).
On \(\operatorname{Int}C_2\), give \(R_{i,j}\) the \(e\)-weight
\[
 \max\{u_{i,j}+c,mb\}.
\]
On \(Y^{\an}\setminus C_{\mathrm{mid}}\), set \(R_{i,j}=mq\).
These are an ordinary open cover, and their prescriptions agree on
the overlap because \(mb>M\) there.
The maximum and locality properties of classical psh functions
\cite[Proposition~4.2]{GR} prove classicality.
The metric difference \(R_{i,j}-mq\) lies in \([0,M]\);
thus it is bounded globally.
Use the same prescription with \(u_i\) to define \(R_i\).
For a decreasing sequence of real numbers,
the maximum with a fixed number commutes with its limit.
It follows that \(R_{i,j}\downarrow R_i\) at every point.
Finally take \(V=\operatorname{Int}C_1\). Here \(b=0\) and the
first branch is positive, proving \eqref{eq:common-collar-weights}.
The closure of \(V\) lies in \(C_1\subset O\).
\end{proof}

\begin{lemma}[One inactive floor]
\label[lemma]{lem:inactive-floor}
In \Cref{lem:common-collar}, suppose that
\(P_{i,j,m}\downarrow R_{i,j}\) are pure Fubini--Study metrics
on \(L\). Their weights on the same \(V\), and all their
nonnegative translates in \(\lambda\mathbb Q\), have finite
tropically convex rational PL presentations in algebraic
coordinates. The algebraic chart and \(V\) are independent of \(i,j,m\).
\end{lemma}
\begin{proof}
After clearing a tensor denominator, each weight in the fixed frame
has the form
\[
 P=\frac1N\max_a\log|s_a/e^{\otimes N}|,
\]
where the finitely many \(s_a\) generate \(NL\).
Rational constants in \(\lambda\mathbb Q\) can be absorbed by
raising to one common tensor power and multiplying sections by
powers of \(\pi\).
All these weights lie above their limits as \(m\to\infty\), and those limits
lie above the finitely many bounded functions \(u_i+c\) on \(V\).
Choose a single $B\in\lambda\mathbb Q$ strictly below all of them.
This floor will keep the tropical expression finite at zeros of
individual sections, while remaining inactive in every approximating
weight.
Then \(P=\max\{B,P\}\) on \(V\), for every \(i,j,m\).

Write \(B=a\lambda/N_0=(1/N_0)\log|\pi^{-a}|\).
Appending the nonzero constant regular function \(\pi^{-a}\)
to the algebraic moment coordinates represents the floor by a
rational linear expression with zero constant term.
In the valuation coordinates \(t_f=-\log|f|\), the resulting function
is a maximum of rational linear functions \(-t_f/N\).
The floor remains finite even when some section coordinates vanish.
Near such a point, those coordinates have value \(+\infty\) and their
branches are inactive; deleting them leaves a finite rational PL
expression on a neighborhood of that point.
Restrict the ambient extended tropical domain to the locus where
the appended constant coordinate is finite. The finitely many
possible subsets of inactive coordinates give the finite PL
cover in Wu's definition. The expression is continuous on this
domain and convex on its finite-coordinate part, hence tropically convex by
\cite[Lemma~3.4(1), Definitions~3.1 and~3.7]{Wu}.
The appended constant coordinate is algebraic; if an affine embedding
is desired, include generators of the fixed affine coordinate ring
as additional coordinates which the displayed function ignores.

A nonnegative translate still lies above \(B\), and a translate in
\(\lambda\mathbb Q\) is handled by the same section-power operation.
Thus the floor and the underlying domain are independent of the indices
and of the translate. Subtracting the fixed \(c\) in
\eqref{eq:common-collar-weights} only replaces the floor by \(B-c\).
This construction does not require setting \(\lambda=1\).
\end{proof}

\begin{lemma}[Ordinary algebraic neighborhoods]
\label[lemma]{lem:local-algebraization}
Let \(X\) be good, Hausdorff, boundaryless and quasi-smooth over \(K\).
For every \(x\in X\) and ordinary neighborhood \(N\) of \(x\),
there is a smaller ordinary neighborhood \(W\subset N\)
which is \(K\)-analytically isomorphic to an ordinary open
of a smooth affine \(K\)-variety.
\end{lemma}
\begin{proof}
Shrink \(N\) inside the interior of an affinoid neighborhood of \(x\);
it is then separated and remains boundaryless.
Ducros's algebraization lemma \cite[Lemma~(0.21)]{Duc}, applied
with the full group of positive radii, gives an affinoid
neighborhood \(V\) of \(x\) in \(N\), identified with an affinoid
domain in the analytification of a smooth affine \(K\)-variety \(Y\).
No extension of \(K\) is made.

The relative-interior composition formula, with locally separated
intermediate spaces \cite[\S2.11(3)--(4)]{GR}, gives
\[
 \operatorname{Int}(V/N)=\operatorname{Int}(V/K)
                      =\operatorname{Int}(V/Y^{\an}),
\]
because \(N\) and \(Y^{\an}\) are boundaryless.
Relative interiors of analytic-domain embeddings are their ordinary
topological interiors; see \cite[Fact~3.4.11 and Fact~4.2.15(i)]{Tem}.
The common interior is therefore ordinary open in both \(N\)
and \(Y^{\an}\), and contains \(x\) since \(V\) was a neighborhood.
It is the required \(W\). One may shrink it further to have compact
closure in a larger such chart.
\end{proof}

\begin{lemma}[Finite constants and integration]
\label[lemma]{lem:finite-constants}
Let \(Y\) be a smooth integral affine \(K\)-variety, and let \(L\)
be the algebraic closure of \(K\) in \(K(Y)\). Then \(L/K\)
is finite separable, \(Y\) is naturally an \(L\)-variety, and
\(Y/L\) is smooth and geometrically integral.
Viewing this same analytic space over \(K\) or over \(L\) does not
change its functions, absolute values or classical psh functions.
If \(d=[L:K]\), the calibrated integration currents satisfy
\begin{equation}\label{eq:finite-constant-calibration}
 [1]_K=d[1]_L.
\end{equation}
Consequently the mixed currents and weighted currents, whenever
constructed by the recursion in the main theorem, satisfy
\[
 T_{k,K}=dT_{k,L},\qquad S_{k,K}=dS_{k,L}.
\]
\end{lemma}
\begin{proof}
The extension \(L/K\) is finite because \(K(Y)/K\) is finitely
generated, and it is separable in characteristic zero.
Every element of \(L\) is integral over \(K\), hence over the normal
affine coordinate ring of \(Y\); it therefore belongs to that ring.
This gives the \(L\)-structure.
The diagonal factor of \(Y\times_K L\), singled out by
\(L\otimes_K L\to L\), is \(Y\) itself as an \(L\)-scheme.
It is open and closed, so it is smooth over \(L\).
The function-field extension \(K(Y)/L\) is regular, giving geometric
integrality.

There is a unique extension of the absolute value of the complete
field \(K\) to \(L\). Thus the two analytifications have the same
points, topology and analytic functions.
An \(E/L\)-curve test maps to the diagonal open-and-closed factor
\(Y_L\times_L E\) of \(Y_K\times_K E\), and is therefore a permitted
\(K\)-curve test. Conversely, for an \(E/K\)-curve test, pass to a
completed algebraic closure \(\Omega/E\). Then
\(L\otimes_K\Omega=\prod_{\sigma:L\to\Omega}\Omega\).
The inverse images of these factors are open-and-closed curves.
Classicality over \(L\) gives subharmonicity on each of them.
Sheafiness and descent under analytic field extension
\cite[Proposition~3.5(7)]{GR} give subharmonicity on the original
\(E\)-curve. Upper semicontinuity and local boundedness
are unchanged.

The moment skeletons are the same in the two interpretations.
Forgetting \(L\) multiplies each finite moment degree by \(d\);
its calibration is therefore multiplied by \(d\), as in
\cite[Remark~14.4.3]{CLD25}.
This proves \eqref{eq:finite-constant-calibration}.
Differentiation and multiplication by a fixed Borel function are
linear. Induction through the recursive definition gives the same
single factor \(d\) at every exterior degree, not \(d^k\).
\end{proof}

\returnlink{prop:ordinary-transport}{transport to ordinary analytic charts}
\returnlink{prop:local-algebraic-regularization}{local algebraic regularization}


\section*{Use of generative AI}
Generative AI was used in developing and drafting mathematical arguments
and in revising the exposition of this manuscript.

\begin{thebibliography}{99}
\raggedright

\bibitem{BMP}
B.~Bhatt, L.~Ma, Z.~Patakfalvi, K.~Schwede, K.~Tucker,
J.~Waldron, and J.~Witaszek, with an appendix by R.~Datta,
\emph{Test ideals in mixed characteristic: a unified theory up to perturbation},
\href{https://arxiv.org/abs/2401.00615v2}{arXiv:2401.00615v2}, 2025.

\bibitem{BM}
E.~Bierstone and P.~D.~Milman,
\emph{Functoriality in resolution of singularities},
Max-Planck-Institut f\"ur Mathematik preprint 2007-30, 2007.
\href{https://archive.mpim-bonn.mpg.de/2624/1/preprint_2007_30.pdf}{Author manuscript}.

\bibitem{BE}
S.~Boucksom and D.~Eriksson,
\emph{Spaces of norms, determinant of cohomology and Fekete points
in non-Archimedean geometry},
\href{https://arxiv.org/abs/1805.01016v3}{arXiv:1805.01016v3}, 2021.

\bibitem{BGM}
S.~Boucksom, W.~Gubler, and F.~Martin,
\emph{Non-Archimedean volumes of metrized nef line bundles},
\'Epijournal de G\'eom\'etrie Alg\'ebrique \textbf{5} (2021), Article 13;
\href{https://arxiv.org/abs/2011.06986v3}{arXiv:2011.06986v3}.

\bibitem{CLD12}
A.~Chambert--Loir and A.~Ducros,
\emph{Formes diff\'erentielles r\'eelles et courants sur les espaces de Berkovich},
\href{https://arxiv.org/abs/1204.6277v1}{arXiv:1204.6277v1}, 2012.
This edition fixes the original test sheaf used here.

\bibitem{CLD25}
A.~Chambert--Loir and A.~Ducros,
\emph{Formes diff\'erentielles r\'eelles et courants sur les espaces de Berkovich},
\href{https://arxiv.org/abs/1204.6277v2}{arXiv:1204.6277v2}, 2025.

\bibitem{Duc}
A.~Ducros,
\emph{Espaces de Berkovich, polytopes, squelettes et th\'eorie des mod\`eles},
Confluentes Mathematici \textbf{4} (2012), no.~4, Article 1250007,
57 pp. The numbering used here is that of the
\href{https://webusers.imj-prg.fr/~antoine.ducros/tropicbis.pdf}{53-page author manuscript},
Lemma~(0.21).

\bibitem{FGK}
Y.~Fang, W.~Gubler, and K.~K\"unnemann,
\emph{On the non-Archimedean Monge--Amp\`ere equation in mixed characteristic},
Nagoya Mathematical Journal \textbf{259} (2025), 548--561;
\href{https://arxiv.org/abs/2203.12282v3}{arXiv:2203.12282v3}.

\bibitem{GJR}
W.~Gubler, P.~Jell, and J.~Rabinoff,
\emph{Forms on Berkovich spaces based on harmonic tropicalizations},
\href{https://arxiv.org/abs/2111.05741v3}{arXiv:2111.05741v3}.

\bibitem{GR}
W.~Gubler and J.~Rabinoff,
\emph{Classically psh and pluriharmonic functions on Berkovich spaces},
\href{https://arxiv.org/abs/2506.13548v2}{arXiv:2506.13548v2}, 2025.

\bibitem{Holm}
H.~Holm,
\emph{The structure of balanced big Cohen--Macaulay modules over
Cohen--Macaulay rings},
Glasgow Mathematical Journal \textbf{59} (2017), 549--561.
\href{https://web.math.ku.dk/~holm/download/MCM-big.pdf}{Author manuscript}.

\bibitem{Stacks}
The Stacks Project Authors,
\emph{The Stacks Project}, 2026.
\href{https://stacks.math.columbia.edu/tag/0BWB}{Section 49.12:
The Noether different};
\href{https://stacks.math.columbia.edu/tag/0G4C}{Section 37.32:
Bertini theorems}.

\bibitem{Tem}
M.~Temkin,
\emph{Introduction to Berkovich analytic spaces},
preliminary lecture notes, 2010.
\href{https://berkovich-2010.imj-prg.fr/lecture-notes/Temkin-lecture-notes.pdf}{Lecture notes}.

\bibitem{Thuillier}
A.~Thuillier,
\emph{Th\'eorie du potentiel sur les courbes en g\'eom\'etrie
analytique non archim\'edienne. Applications \`a la th\'eorie d'Arakelov},
doctoral thesis, Universit\'e de Rennes~1, 2005.
\href{https://theses.hal.science/tel-00010990/document}{Thesis}.

\bibitem{Wanner}
V.~Wanner,
\emph{Comparison of two notions of subharmonicity on non-archimedean curves},
\href{https://arxiv.org/abs/1801.04713v3}{arXiv:1801.04713v3}, 2018.

\bibitem{Wu}
L.~Wu,
\emph{On the local theory of non-archimedean psh functions},
\href{https://arxiv.org/abs/2609.27610v1}{arXiv:2609.27610v1}, 2026.

\end{thebibliography}

\clearpage
\section*{Index of notation}\label{sec:notation}
The entries are grouped by mathematical role. Each locator links
to the definition or principal use and gives its page.

\small
\begin{longtable}{@{}>{\raggedright\arraybackslash}p{.18\textwidth}
>{\raggedright\arraybackslash}p{.48\textwidth}
>{\raggedright\arraybackslash}p{.26\textwidth}@{}}
\toprule
Symbol & Meaning and scope & Definition or first use\\
\midrule
\endfirsthead
\toprule
Symbol & Meaning and scope & Definition or first use\\
\midrule
\endhead
\midrule
\multicolumn{3}{r}{\emph{Continued on the next page}}\\
\endfoot
\bottomrule
\endlastfoot

\multicolumn{3}{@{}l}{\textbf{Fields, domains and input classes}}\\*[3pt]
\(K,\pi,\lambda\) &
Finite extension of \(\mathbb Q_p\), chosen uniformizer, and
\(\lambda=-\log|\pi|>0\) for the prescribed absolute value. &
\Cref{conv:spaces}, p.~\pageref{conv:spaces}\\
\(X,U,n\) &
Good Hausdorff boundaryless quasi-smooth analytic space, ordinary
open, and pure dimension. &
\Cref{conv:spaces}, p.~\pageref{conv:spaces}\\
\(\Cb(U)\) &
Locally bounded upper-semicontinuous classical psh functions;
curve tests after every complete isometric field extension. &
\Cref{def:classes}, p.~\pageref{def:classes}\\
\(\PLsp(U)\) &
Continuous model-semipositive PL functions. Rational slopes
and constants in \(\lambda\mathbb Q\) are used in the constructions. &
\Cref{def:classes}, p.~\pageref{def:classes}\\
\(\Rb(U)\) &
Locally bounded functions locally obtained as decreasing
PL limits on one fixed ordinary domain. &
\Cref{def:classes}, p.~\pageref{def:classes}\\
\(\mathcal G_K\) &
Fields obtained from \(K\) by finitely many unit-Gauss completions;
the simultaneous dimension induction uses arbitrary data over them. &
\Cref{def:gauss-fields,lem:gauss-countable},
p.~\pageref{lem:gauss-countable}\\
\(L/K,\ d\) &
Finite algebraic constant field of a smooth integral affine
\(K\)-variety, and its degree; forgetting constants multiplies
each calibrated current by the single factor \(d\). &
\Cref{lem:finite-constants}, p.~\pageref{lem:finite-constants}\\[4pt]

\multicolumn{3}{@{}l}{\textbf{Forms, coefficient measures and mixed currents}}\\*[3pt]
\(A^{p,q}, A_c^{p,q}\) &
Original CLD smooth forms and compact test forms, with their
canonical calibration. &
\Cref{conv:forms}, p.~\pageref{conv:forms}\\
\(D,\dpr,\dpp\) &
The original differentials and curvature operator
\(D=\dpr\dpp\). &
\Cref{conv:forms}, p.~\pageref{conv:forms}\\
\([1]\) &
Integration current in the ground-field calibration. &
\eqref{eq:recursive-current}, p.~\pageref{eq:recursive-current}\\
\(\mu_{T,\eta},hT\) &
Radon measure on the analytic domain representing the scalar functional
\(a\mapsto\langle T,a\eta\rangle\), and multiplication by a
locally bounded Borel function. &
\Cref{def:coefficient-radon}, p.~\pageref{def:coefficient-radon}\\
\(T_k(u_1,\ldots,u_k)\) &
Symmetric closed strongly positive mixed curvature current of
bidegree \((k,k)\). &
\Cref{thm:mixed-currents}, p.~\pageref{thm:mixed-currents}\\
\(S_k(u_0;u_1,\ldots,u_k)\) &
Weighted current \(u_0T_k(u_1,\ldots,u_k)\), with undifferentiated
argument \(u_0\). &
\eqref{eq:source-current}, p.~\pageref{eq:source-current}\\
\(W_r,\ A_r\) &
Span of decomposable squares in the symmetric tensor space
of \(r\)-vectors, and its annihilator in the dual tensor space. &
\Cref{lem:strong-coefficients}, p.~\pageref{lem:strong-coefficients}\\
\(\MA_{\rm PL}\) &
Top-degree measure in the local PL theory, compared with the
continuous CLD product. &
\Cref{prop:wu-input}, p.~\pageref{prop:wu-input}\\
\(j,\ |\cdot|_*\) &
Transport under \(|\cdot|_*=|\cdot|^{1/\lambda}\), used only
to compare valuation conventions with the PL input. &
\eqref{eq:valuation-calibration}, p.~\pageref{eq:valuation-calibration}\\[4pt]

\multicolumn{3}{@{}l}{\textbf{Metrics, energy and slicing}}\\*[3pt]
\(c_1(L,\phi)\) &
Curvature of the additive metric \(\phi\), with
model mass equal to structural degree. &
\Cref{conv:forms}, \eqref{eq:native-model-measure},
p.~\pageref{eq:native-model-measure}\\
\(\Gamma(g),\Gamma(g,h)\) &
First-order quadratic energy current and its symmetric
polarization. &
\eqref{eq:gamma}, p.~\pageref{eq:gamma}\\
\(E_\Omega(g,h)\) &
Energy paired with a smooth strongly positive closed
complementary background. &
\Cref{prop:first-order-energy}, p.~\pageref{prop:first-order-energy}\\
\(P_L\rho\) &
Pointwise envelope of continuous model-semipositive minorants
of \(\rho\) on the same ample line \(L\). &
\Cref{def:envelope}, p.~\pageref{def:envelope}\\
\(E,\ H,\ \iota\) &
Unit-Gauss parameter field, its generic hyperplane, and the map
from that hyperplane to the original projective analytic space. &
\Cref{lem:single-hyperplane}, p.~\pageref{lem:single-hyperplane}\\
\(\Delta_C b,\Gamma_C\) &
Positive curve Laplacian of a subharmonic function,
and the curve energy measure. &
\Cref{lem:curve-native,lem:bounded-curve-energy},
p.~\pageref{lem:bounded-curve-energy}\\
\(p:W\to\mathbb G_m^{n-1},\ \Omega\) &
Algebraic curve-slicing map and a nonnegative top form in
the specified standard torus coordinates. &
\Cref{lem:borel-coarea}, p.~\pageref{lem:borel-coarea}\\
\(\nu_\Omega,\ \mu_y\) &
Base skeletal measure and the curve Laplacian kernel,
defined off one common null parameter set. &
\Cref{thm:actual-kernel}, p.~\pageref{thm:actual-kernel}\\[4pt]

\multicolumn{3}{@{}l}{\textbf{Normal-model envelopes and local geometry}}\\*[3pt]
\(R,F\) &
In Appendix~\ref{sec:normal-envelopes}, a complete
mixed-characteristic DVR and its characteristic-zero
fraction field; no residue-field finiteness is assumed. &
\Cref{sec:normal-envelopes}, p.~\pageref{sec:normal-envelopes}\\
\(S,W,E,F_S\) &
Complete local normal domain, its chosen canonical module
\(W=\cO_S(-\Delta)\), and their top local cohomology modules. &
\Cref{sec:normal-envelopes}, p.~\pageref{sec:normal-envelopes}\\
\(\mathfrak N_{S/A}\) &
Noether different of a finite extension from a regular local
subring; its witnesses give the flat factorization. &
\Cref{lem:different-factor}, p.~\pageref{lem:different-factor}\\
\(\tau(\cO_{\mathcal X},\Delta,\mathfrak a^s)\) &
Mixed-characteristic test ideal; the global fixed correction
is independent of the ideals and rational exponents. &
\Cref{thm:fixed-correction}, p.~\pageref{thm:fixed-correction}\\
\(\mathfrak c_\ell,\ \beta_m,\ \psi_m\) &
Weighted base ideals, stable asymptotic test ideals, and their
normalized logarithmic norms in the envelope construction. &
\eqref{eq:weighted-ideals}, p.~\pageref{eq:weighted-ideals}\\
\(C_1,C_{\rm mid},C_2,\ b\) &
Nested compact norm boxes and their logarithmic collar barrier. &
\Cref{lem:common-collar}, p.~\pageref{lem:common-collar}\\
\(R_{i,j},P_{i,j,m},Q_{i,j}\) &
Bounded projective metrics on a common line bundle, decreasing pure-FS
rows, and the compact triangular diagonal. &
\Cref{lem:common-collar,lem:triangular-diagonal},
p.~\pageref{lem:triangular-diagonal}\\
\end{longtable}
\normalsize

\end{document}
